% !TeX root = ../../Book5.tex
% 纲要标题：## 第21章*　Hopf 代数与仿射代数群
% 纲要位置：工作记录/计划纲要.md，第 4755 行。

\OptionalChapter{Hopf 代数与仿射代数群}{b5:ch:21}
\BookChapterNumber{b5:ch:21}

\begin{chapterabstract}
	本章建立 Hopf 代数的乘法、余乘法、余单位和对极，比较群代数、包络代数与坐标环，并用余模和矩阵系数研究有理表示与群概形例子。
\end{chapterabstract}

\begin{learninggoals}
	\begin{itemize}
		\item 在具体生成元上核验余乘法、余单位和对极，并构造张量与对偶作用。
		\item 区分群代数和坐标代数、模与余模、基域点与所有交换代数上的点。
		\item 从余模证明有理表示的局部有限性、环面权分解及仿射群的线性实现。
		\item 说明仅由微分表示无法恢复哪些整数与正特征信息，比较等距、相似和对合自同构群。
	\end{itemize}
\end{learninggoals}

乘法群的元素可以相乘，也可以取逆。若把它看成由坐标函数 \(t\) 与 \(t^{-1}\) 描述的对象，两个群元素相乘时，\(t\) 的取值随之相乘；在坐标环中，这表现为 \(\Delta(t)=t\otimes t\)。单位元给出 \(\varepsilon(t)=1\)，取逆则给出 \(S(t)=t^{-1}\)。原来作用在点上的三种运算，如今变成作用在函数上的代数映射，并满足相应的恒等式。沿着这个具体例子，Hopf 代数的公理和表示张量积所需的余乘法就有了可核对的来源。

\ProseParagraphBreak
本章使用张量积、代数同态和第三卷的仿射坐标环语言；群概形从可表示函子及坐标方程出发介绍。中心商的局部计算还使用第三卷第24.3节的射影仿射图和第12.5节的忠实平坦性，平展性判别使用第三卷第14章的微分展示。包络代数例子调用本卷第15章，四元数例子则把所需乘法与范数公式重新列出。

\ProseParagraphBreak
Hopf 代数的结构及对偶可参阅\cite[Sections 5.2--5.3]{EtingofEtAl2015TensorCategories}。该书的 Hopf 定义要求对极可逆，本章采用更一般的定义并在使用对偶时说明实际所需条件；所讨论的群代数、包络代数和交换坐标代数不受这一区别影响。

% !TeX root = ../../Book5.tex
% 纲要标题：### 21.1 余代数与 Hopf 代数
% 纲要位置：工作记录/计划纲要.md，第 4757 行。

\section{余代数与 Hopf 代数}
\label{b5:sec:2101}


% 纲要细目（保留原定写作范围）：
%
% * 余乘法、余单位与对极
% * 群代数及坐标环的对偶方向
% * 模与余模的区别
% * 证明有限群的四种模与余模对应，展开 Casimir 的张量交叉项，并以求值配对解释逆转置
%

群表示与 Lie 代数表示都能作张量积，但作用公式不同：群元素同时作用于两个因子，Lie 代数元素则分别作用后相加。余乘法把这两种规则统一写成代数中的映射；再把单位表示和对偶表示所需的规则一并加入，就得到 Hopf 代数。

\ProseParagraphBreak

前面用 \(k[G]\) 表示群代数。本章将群代数写作 \(kG\)，有限群上的函数代数写作 \(k^G\)，而将 \(k[G]\) 留给仿射群概形的坐标环。这些对象的乘法和余乘法分别定义，不能仅凭相同的记号互相识别。

\subsection{余乘法、余单位与对极}
\label{b5:subsec:2101:01}

\begin{definition}\label{b5:def:coalgebra-hopf}
	设 \(k\) 为域，\(H\) 为 \(k\) 向量空间。线性映射 \(\Delta:H\to H\otimes_kH\)、\(\varepsilon:H\to k\) 满足
	\[
		(\Delta\otimes1)\Delta=(1\otimes\Delta)\Delta,\qquad (\varepsilon\otimes1)\Delta=1_H=(1\otimes\varepsilon)\Delta
	\]
	时，称 \((H,\Delta,\varepsilon)\) 为\term[yudaishu]{余代数}[coalgebra]，称两映射为余乘法与余单位；后两式用 \(k\otimes H=H=H\otimes k\) 识别。 若 \(H\) 还为含幺结合 \(k\) 代数，且 \(\Delta,\varepsilon\) 是含幺代数同态，就称为\term[shuangdaishu]{双代数}[bialgebra]。若另有线性映射 \(S:H\to H\) 满足
	\[
		m(S\otimes1)\Delta=\eta\varepsilon=m(1\otimes S)\Delta,
	\]
	其中 \(m\) 是乘法、\(\eta:k\to H\) 是单位映射，就称为\term[Hopf-daishu]{Hopf 代数}[Hopf algebra]，称 \(S\) 为\term[duiji]{对极}[antipode]。
\end{definition}
\symidx{\(\Delta,\varepsilon,S\)：Hopf 代数的余乘法、余单位与对极}
上述定义允许一般对极；本章的群代数、包络代数及交换坐标代数例子还具有可逆对极。Etingof 等的《Tensor Categories》在定义 Hopf 代数时要求对极可逆，因此引用其张量范畴结论时须保留这一差别。
\ProseParagraphBreak
常把 \(\Delta(h)\) 写为 \(\sum h_{(1)}\otimes h_{(2)}\)，称为 Sweedler 记号。它表示一个有限和，不表示每个 \(h\) 都只有一个纯张量分解。余结合律允许把三重余乘法统一写为 \(\sum h_{(1)}\otimes h_{(2)}\otimes h_{(3)}\)。
\symidx{\(\sum h_{(1)}\otimes h_{(2)}\)：余乘法的 Sweedler 记号}
\begin{proposition}\label{b5:prop:antipode-unique}
	设 \(k\) 为域，\(H\) 为 \(k\) 双代数。其对极 \(S:H\to H\) 若存在，则唯一，并满足 \(S(1)=1\)、\(S(ab)=S(b)S(a)\) 对全部 \(a,b\in H\) 成立。
\end{proposition}
\begin{proof}
	在 \(\operatorname{Hom}_k(H,H)\) 上定义卷积 \((f*g)(h)=\sum f\bigl(h_{(1)}\bigr)g\bigl(h_{(2)}\bigr)\)。 结合律来自 \(m\) 的结合与 \(\Delta\) 的余结合，单位为 \(\eta\varepsilon\)。对极就是恒等映射的双边卷积逆，故唯一；代入单位即得 \(S(1)=1\)。 对反乘法性质，改在 \(\operatorname{Hom}_k(H\otimes H,H)\) 中卷积，张量积余乘法把 \(a\otimes b\) 送到 \(\sum\bigl(a_{(1)}\otimes b_{(1)}\bigr)\otimes\bigl(a_{(2)}\otimes b_{(2)}\bigr)\)。 这里的卷积单位为 \((a,b)\mapsto\varepsilon(a)\varepsilon(b)1\)。由 \(\Delta(ab)=\sum a_{(1)}b_{(1)}\otimes a_{(2)}b_{(2)}\)，对极等式给出 \((S\circ m)*m(a,b)=\varepsilon(ab)1=\varepsilon(a)\varepsilon(b)1\)，另一侧同理，所以 \(S\circ m\) 是乘法映射 \(m\) 的卷积逆。另一个映射 \(F(a,b)=S(b)S(a)\) 也给出逆：两侧复合分别为
	\[
		\sum S\bigl(b_{(1)}\bigr)S\bigl(a_{(1)}\bigr)a_{(2)}b_{(2)} =\varepsilon(a)\varepsilon(b)1
	\]
	（先用 \(a\) 的对极等式，再用 \(b\) 的），另一侧先消去 \(b\)，得到 \(\sum a_{(1)}b_{(1)}S\bigl(b_{(2)}\bigr)S\bigl(a_{(2)}\bigr)=\varepsilon(a)\varepsilon(b)1\)。 逆的唯一性使两映射相同。
\end{proof}
记 \(\tau(a\otimes b)=b\otimes a\) 为张量因子翻转。若 \(\tau\Delta=\Delta\)，则称余交换。乘法交换与余乘法余交换是两项不同条件，以下两种来自同一个群的代数将说明这种区别。

\subsection{群代数及坐标环的对偶方向}
\label{b5:subsec:2101:02}

\begin{proposition}\label{b5:prop:group-enveloping-hopf}
	设 \(k\) 为域，\(G\) 为群，\(\mathfrak g\) 为 \(k\) 上 Lie 代数。下列公式分别给出 \(kG\) 与 \(U(\mathfrak g)\) 上的 Hopf 代数结构：
	\[
		\begin{array}{c|ccc}
			&H:\Delta&\varepsilon&S\\ \hline kG&g\mapsto g\otimes g&g\mapsto1&g\mapsto g^{-1}\\
			U(\mathfrak g)&x\mapsto x\otimes1+1\otimes x&x\mapsto0&x\mapsto-x
		\end{array}
	\]
	第二行的公式在 \(x\in\mathfrak g\) 上规定，再按相应代数或反代数映射延拓。两者都余交换。
\end{proposition}
\begin{proof}
	群代数的各式在群基元上核验公理即可，例如 \(m(S\otimes1)\Delta(g)=g^{-1}g=1=\varepsilon(g)1\)。 包络代数中，映射 \(x\mapsto x\otimes1+1\otimes x\) 保持 Lie 括号，因为交叉的两个因子交换，故由泛性质延拓为代数同态。增广零化全部 Lie 括号；映射 \(x\mapsto-x\) 到反代数保持 Lie 括号，因而给出反同态 \(S\)。余结合及余单位等式在生成元上成立，所以在全部代数上成立。对极等式可对字长归纳：若两较短字 \(u,v\) 满足左式，则
	\[
		\sum S\bigl((uv)_{(1)}\bigr)(uv)_{(2)} =\sum S\bigl(v_{(1)}\bigr)S\bigl(u_{(1)}\bigr)u_{(2)}v_{(2)} =\varepsilon(u)\varepsilon(v)1.
	\]
	右式同理；线性性完成证明。
\end{proof}
群代数作为普通代数保存的资料少于它作为 Hopf 代数保存的资料。第2章中，\(\mathbb C[D_8]\) 与 \(\mathbb C[Q_8]\) 都同构于 \(\mathbb C^4\times M_2(\mathbb C)\)，所以仅凭其模范畴不能区分这两个群。余乘法却记住了群基元：若 \(a=\sum_g a_gg\) 满足 \(\Delta(a)=a\otimes a\)、\(\varepsilon(a)=1\)，比较 \(g\otimes h\) 的系数得到
\[
	a_ga_h=0\quad(g\ne h),\qquad a_g^2=a_g, \qquad \sum_g a_g=1.
\]
在域中只能有一个系数非零，且它等于一，故这样的 \(a\) 恰为原群元素。Hopf 同构保持这两项等式及乘法，因而恢复群同构；上述普通代数同构不能自动保持 Hopf 结构。余乘法与余单位也正是下文规定表示的张量积与单位对象的资料。

若 \(G\) 有限，函数空间 \(k^G\) 则按逐点乘法成为交换 Hopf 代数：
\[
	\Delta f(g,h)=f(gh),\qquad \varepsilon(f)=f(1),\qquad Sf(g)=f(g^{-1}).
\]
\symidx{\(k^G\)：有限群的函数 Hopf 代数}
在点函数基 \(\delta_x\) 中， \(\Delta\delta_x=\sum_{ab=x}\delta_a\otimes\delta_b\)。 这是 \(kG\) 的线性对偶 Hopf 代数；函数代数通常不余交换，恰在 \(G\) 交换时余交换。群无限时，任意函数的两变量拉回未必在代数张量积 \(k^G\otimes k^G\) 中，因此不能不加限制照搬有限对偶。

\ProseParagraphBreak
例如在加法群 \(\mathbb Z\) 上取点函数 \(\delta_0\)。乘法拉回在 \((m,n)\) 处为 \(\delta_{m+n,0}\)。若它来自代数张量积，就可写成某个有限和
\[
	\delta_{m+n,0}=\sum_{i=1}^r a_i(m)b_i(n).
\]
固定 \(m\) 后，这个关于 \(n\) 的函数必落在 \(b_1,\ldots,b_r\) 的线性生成空间中。但左边随 \(m\) 变化给出所有点函数 \(\delta_{-m}\)，它们线性无关：在各自支集点求值即可逐个读出系数。这与有限维生成空间矛盾，故该两变量函数确实不在 \(k^{\mathbb Z}\otimes k^{\mathbb Z}\) 中。

\subsection{模与余模的区别}
\label{b5:subsec:2101:03}

有限群表示在基下写成 \(gv_j=\sum_i a_{ij}(g)v_i\)。与其对每个 \(g\) 分别保存这个矩阵，也可以把系数看成函数 \(a_{ij}\in k^G\)，一次写成 \(v_j\mapsto\sum_i v_i\otimes a_{ij}\)。群乘法在系数函数上的拉回，正好给出这种展开应满足的相容条件。这便是余模的来源。

\begin{definition}\label{b5:def:right-comodule}
	设 \(k\) 为域，\((C,\Delta,\varepsilon)\) 为 \(k\) 余代数。\(k\) 向量空间 \(V\) 配 \(k\) 线性映射 \(\delta:V\to V\otimes_k C\)，使
	\[
		(\delta\otimes1)\delta=(1\otimes\Delta)\delta,\qquad (1\otimes\varepsilon)\delta=1_V,
	\]
	时，称为右 \(C\)\term[yumo]{余模}[comodule]。第一式将两种继续展开余作用的次序联系起来，即下图交换：
	\[
		\begin{tikzcd}[column sep=large,row sep=large]
			V\arrow[r,"\delta"]\arrow[d,"\delta"'] &V\otimes_kC\arrow[d,"1_V\otimes\Delta"]\\
			V\otimes_kC\arrow[r,"\delta\otimes1_C"'] &V\otimes_kC\otimes_kC
		\end{tikzcd}
	\]
	给定两个右 \(C\) 余模 \((V,\delta_V)\)、\((W,\delta_W)\)，若 \(k\) 线性映射 \(f:V\to W\) 使下图交换，就称它为余模同态：
	\[
		\begin{tikzcd}[column sep=large,row sep=large]
			V\arrow[r,"\delta_V"]\arrow[d,"f"'] &V\otimes_kC\arrow[d,"f\otimes1_C"]\\
			W\arrow[r,"\delta_W"']&W\otimes_kC
		\end{tikzcd}
	\]
\end{definition}
\symidx{\(\delta:V\to V\otimes_kC\)：右余模的余作用}
模是代数元素作用于向量的规则；余模是把一个向量展开成向量与系数函数的规则。只有适当的对偶条件成立时，两种语言才可直接转换。例如有限维代数 \(A\) 的左模等同于右 \(A^*\) 余模；取对偶基 \(a_i,a^i\)，余作用为 \(v\mapsto\sum_i a_iv\otimes a^i\)。 代数乘法公理与单位公理分别对应余结合与余单位。

\begin{proposition}\label{b5:prop:finite-group-module-comodule}
	设 \(k\) 为域，\(G\) 为有限群，并以 \(k^G=(kG)^*\) 表示函数 Hopf 代数。左 \(kG\) 模和右 \(k^G\) 余模都与 \(G\) 在 \(k\) 向量空间上的线性表示对应；左 \(k^G\) 模和右 \(kG\) 余模都与 \(G\)-分次向量空间对应。这些对应保留底层向量空间、相应的同态及直和分解。
\end{proposition}
\begin{proof}
	左 \(kG\) 模的作用限制到群基 \(G\)，便是群表示；反过来把群作用线性延拓即可。对这样的作用，令
	\[
		\delta(v)=\sum_{g\in G}gv\otimes\delta_g.
	\]
	由 \(g(hv)=(gh)v\) 和 \(\Delta\delta_x=\sum_{ab=x}\delta_a\otimes\delta_b\)，余结合律两边都为 \(\sum_{g,h}ghv\otimes\delta_g\otimes\delta_h\)，而 \(1v=v\) 给出余单位律。反过来，对右 \(k^G\) 余模取 \(gv=(1\otimes\operatorname{ev}_g)\delta(v)\)，其中 \(\operatorname{ev}_g(f)=f(g)\)，同样的两个等式分别给出群乘法与单位作用。这两项构造互为逆。

	函数代数中的 \(\delta_g\) 是和为单位的正交幂等元，所以每个左 \(k^G\) 模都有分解
	\[
		V=\bigoplus_{g\in G}V_g,\qquad V_g=\delta_gV.
	\]
	给定这样的分次，令 \(f\in k^G\) 在 \(V_g\) 上按标量 \(f(g)\) 作用，就恢复原来的模。另一方面，分次空间上的公式 \(\delta(v)=v\otimes g\)、\(v\in V_g\)，给出右 \(kG\) 余模，因为 \(\Delta(g)=g\otimes g\)、\(\varepsilon(g)=1\)。任一右 \(kG\) 余作用可唯一写为 \(\delta(v)=\sum_g p_g(v)\otimes g\)。比较余结合律中 \(a\otimes b\) 的系数，得到
	\[
		p_ap_b=0\quad(a\ne b),\qquad p_g^2=p_g;
	\]
	余单位律又给出 \(\sum_g p_g=1_V\)。因此 \(V=\bigoplus_g\operatorname{im}p_g\)，且在 \(\operatorname{im}p_g\) 上余作用为 \(v\mapsto v\otimes g\)。上述各式也说明，同态条件分别是与群作用相容或保持各分次片，所以这些对应保留同态与直和。
\end{proof}
可见，将有限维代数取线性对偶，会把它的左模转换为对偶余代数的右余模；它没有把左模直接转换为对偶代数的左模。群表示与分次空间因而对应两种不同的范畴，不能只因 \(kG\) 与 \(k^G\) 互为对偶就将它们认同。

\ProseParagraphBreak
这个区别在模特征中尤其明显。取 \(k=\mathbb F_2\)、\(G=C_2=\{1,s\}\)，令 \(u=s-1\)。由\cref{b5:thm:cyclic-p-modules}在 \(p=2,r=1\) 时的结论，群代数为 \(k[u]/(u^2)\)，其二维不可分解模在基 \(v_1,v_2\) 下满足
\[
	sv_1=v_1,\qquad sv_2=v_2+v_1,\qquad
	[s]_V=\begin{pmatrix}1&1\\0&1\end{pmatrix}.
\]
若它能分成两条稳定直线，\(s\) 在每条直线上的标量都满足 \(\lambda^2=1\)，在 \(\mathbb F_2\) 中只能为一，从而 \(s\) 在整个空间上恒等，与所列矩阵矛盾。相应的右 \(k^G\) 余作用为
\[
	\begin{aligned}
		\delta(v_1)&=v_1\otimes(\delta_1+\delta_s),\\
		\delta(v_2)&=v_2\otimes(\delta_1+\delta_s)+v_1\otimes\delta_s.
	\end{aligned}
\]
由于对应保留直和，这个余模也不可分解。与此相对，函数代数 \(k^G\cong k\times k\) 的左模总有 \(V=\delta_1V\oplus\delta_sV\)；在每个分次片上，任一线性子空间都是子模，有限维时可继续分成一维子模。因此这里的非分裂二维群表示属于右 \(k^G\) 余模范畴，而不属于左 \(k^G\) 模范畴。

\ProseParagraphBreak
余作用还可以把一个多项式参数记录在系数中。设 \(k\) 为域，\(C=k[t]\) 的余乘法为 \(\Delta(t)=t\otimes1+1\otimes t\)，余单位为 \(\varepsilon(t)=0\)。在 \(V=kv_1\oplus kv_2\) 上令
\[
	\delta(v_1)=v_1\otimes1,\qquad \delta(v_2)=v_2\otimes1+v_1\otimes t.
\]
在 \(v_2\) 上，余结合律两边都等于 \(v_2\otimes1\otimes1+v_1\otimes t\otimes1+v_1\otimes1\otimes t\)，余单位律也直接成立。把 \(t\) 代成 \(c\in k\)，便得到矩阵
\(\bigl(\begin{smallmatrix}1&c\\0&1\end{smallmatrix}\bigr)\)；余乘法中的两个加项对应参数相加。直线 \(kv_1\) 是子余模，商余模平凡，但将 \(v_2\) 换成 \(v_2+av_1\) 仍留下 \(v_1\otimes t\) 项，故这条子余模没有补。

\begin{proposition}\label{b5:prop:hopf-tensor-dual}
	设 \(k\) 为域，\((H,\Delta,\varepsilon,S)\) 为 \(k\) 上的 Hopf 代数，\(V,W\) 为幺性左 \(H\) 模。对 \(h\in H\)，用 Sweedler 记号写 \(\Delta(h)=\sum h_{(1)}\otimes h_{(2)}\)。则张量积 \(V\otimes_kW\) 通过
	\[
		h(v\otimes w)=\sum h_{(1)}v\otimes h_{(2)}w
	\]
	成为左模，\(k\) 通过 \(\varepsilon\) 成为单位模。若 \(V\) 有限维，则
	\[
		(h\varphi)(v)=\varphi\bigl(S(h)v\bigr)
	\]
	给出 \(V^*\) 的左模结构。对 \(V\) 的基 \(e_1,\ldots,e_r\) 及对偶基 \(e_1^*,\ldots,e_r^*\)，求值与余求值映射
	\[
		\operatorname{ev}:V^*\otimes_kV\to k,\qquad \varphi\otimes v\mapsto\varphi(v),
	\]
	\[
		\operatorname{coev}:k\to V\otimes_kV^*,\qquad 1\mapsto\sum_i e_i\otimes e_i^*
	\]
	均为 \(H\) 模同态，并使 \(V^*\) 成为 \(V\) 的左对偶。
\end{proposition}
\begin{proof}
	\(\Delta\) 保持乘法使两次作用等于乘积作用，余结合使三重张量的两种括号相容，余单位使 \(k\otimes V\to V\) 相容。对偶上依次作用 \(a,b\) 得到 \(\varphi\bigl(S(b)S(a)v\bigr)=\varphi\bigl(S(ab)v\bigr)\)，所以是左作用。求值映射的等变性为
	\[
		\sum\bigl(h_{(1)}\varphi\bigr)\bigl(h_{(2)}v\bigr) =\varphi\left(\sum S\bigl(h_{(1)}\bigr)h_{(2)}v\right) =\varepsilon(h)\varphi(v).
	\]
	为核验余求值，使用典范同构
	\[
		\Theta:V\otimes_kV^*\xrightarrow{\sim}\operatorname{End}_k(V),\qquad \Theta(v\otimes\varphi)(x)=\varphi(x)v.
	\]
	张量 \(\operatorname{coev}(1)\) 对应 \(1_V\)，因而不依赖基的选择。对任意 \(x\in V\)，有
	\[
		\begin{aligned}
			\Theta\bigl(h\operatorname{coev}(1)\bigr)(x) &=\sum_{i,(h)} e_i^*\bigl(S(h_{(2)})x\bigr)h_{(1)}e_i\\
			&=\sum_{(h)}h_{(1)}S(h_{(2)})x =\varepsilon(h)x.
		\end{aligned}
	\]
	由 \(\Theta\) 单射可知 \(h\operatorname{coev}(1)=\varepsilon(h)\operatorname{coev}(1)\)，故余求值也是 \(H\) 模同态。略去通常的结合与单位同构，两条蛇形复合分别将
	\[
		v\longmapsto\sum_i e_i\,e_i^*(v)=v,\qquad \varphi\longmapsto\sum_i\varphi(e_i)e_i^*=\varphi.
	\]
	这就给出左对偶的两条恒等式；以上计算只用了对极公理。
\end{proof}
在群代数上，张量公式就是对角群作用；在包络代数上，它是 Lie 代数的 Leibniz 型作用。

\ProseParagraphBreak
有限群的分次也与张量作用相容。若 \(V=\bigoplus_aV_a\)、\(W=\bigoplus_bW_b\) 对应左 \(k^G\) 模，\(v\in V_a\)、\(w\in W_b\)，则由 \(\Delta\delta_g=\sum_{xy=g}\delta_x\otimes\delta_y\)，\(\delta_g\) 在 \(v\otimes w\) 上的作用仅在 \(g=ab\) 时为恒等。因此
\[
	(V\otimes_kW)_g=\bigoplus_{ab=g}V_a\otimes_kW_b.
\]
在相应的右 \(kG\) 余模中，同一规则写成 \(\delta(v\otimes w)=v\otimes w\otimes ab\)。这里相乘的是分次指标；右 \(k^G\) 余模所对应的群表示则按对角方式作用于张量积。

\ProseParagraphBreak
包络代数中，生成元的本原性并不延伸到任意中心元素。
\begin{proposition}\label{b5:prop:casimir-coproduct}
	在 \(U\bigl(\mathfrak{sl}_2(\mathbb C)\bigr)\) 的 Hopf 结构中，取标准基 \(e=E_{12}\)、\(f=E_{21}\)、\(h=E_{11}-E_{22}\)，并令 \(\Omega=h^2+2h+4fe\)。则
	\[
		\Delta(\Omega)=\Omega\otimes1+1\otimes\Omega +2h\otimes h+4f\otimes e+4e\otimes f.
	\]
\end{proposition}
\begin{proof}
	余乘法保持乘法，且 \(e,f,h\) 都是本原元，所以
	\[
		\begin{aligned}
			\Delta(h^2)&=h^2\otimes1+1\otimes h^2+2h\otimes h,\\
			\Delta(fe)&=fe\otimes1+1\otimes fe+f\otimes e+e\otimes f.
		\end{aligned}
	\]
	再加上 \(2\Delta(h)\)，便得到所示等式。
\end{proof}
\symidx{\(\Omega=h^2+2h+4fe\)：\(\mathfrak{sl}_2\) 的 Casimir 元的八倍归一化}
\cref{b5:prop:sl2-casimir-central}中的 \(C=h^2+2h+4fe\) 就是这里的 \(\Omega\)，故它属于中心。若以 \(C_K\) 表示第20.2节按 Killing 型归一化的 Casimir 元，则 \(C_K=\Omega/8\)；同一个元素在第14—15章采用了清除分母的归一化。这些交叉项来自乘积展开，不能由中心性消去。

\ProseParagraphBreak
交叉项在最小的张量积中已会改变谱。由\cref{b5:lem:sl2-casimir-operator}，\(\Omega\) 在标准不可约模 \(L_m\) 上按 \(m(m+2)\) 作用，所以在 \(L_1\) 上为 \(3I\)。若只保留 \(\Omega\otimes1+1\otimes\Omega\)，它在 \(L_1\otimes L_1\) 上便是 \(6I\)。取自然二维模的基 \(v_+,v_-\)，使
\[
	hv_+=v_+,\quad hv_-=-v_-,\quad ev_-=v_+,\quad fv_+=v_-,\quad ev_+=fv_-=0.
\]
记交叉项的作用算子为 \(Q=2h\otimes h+4f\otimes e+4e\otimes f\)。它在 \(v_+\otimes v_+\)、\(v_-\otimes v_-\) 上都为标量二；在有序基 \(v_+\otimes v_-,v_-\otimes v_+\) 下则为
\[
	[Q]=\begin{pmatrix}-2&4\\4&-2\end{pmatrix},\qquad
	[\Omega]=6I+[Q]=\begin{pmatrix}4&4\\4&4\end{pmatrix}.
\]
因而中间两项的和、差分别具有 \(\Omega\) 特征值八、零，而两端基张量的特征值都是八。\cref{b5:prop:sl2-tensor-square-projections}已将三维对称子空间识别为 \(L_2\)，将一维交替子空间识别为 \(L_0\)；上述余乘法计算正好给出
\[
	\Omega|_{L_2}=8I,\qquad \Omega|_{L_0}=0.
\]
交叉项中的 \(e\otimes f\) 与 \(f\otimes e\) 把中间两条基张量互相联系，使其对称和与交替差得到不同特征值。这也恢复了第14章的投影公式：若 \(\tau\) 交换两个因子，则 \(\Omega=4(I+\tau)\)，从而 \(\dfrac{\Omega}{8}=\dfrac{I+\tau}{2}\) 与 \(I-\dfrac{\Omega}{8}=\dfrac{I-\tau}{2}\) 分别投到 \(L_2\) 和 \(L_0\)。中心性保证它在这些不可约分量上按标量作用，却不保证张量积中的标量等于两个因子的标量之和。

\ProseParagraphBreak
对偶公式中的对极用来保持求值配对。若群元素 \(g\) 在有限维表示 \(V\) 上的作用矩阵为 \(\rho(g)\)，则在对偶基下，\(S(g)=g^{-1}\) 给出的矩阵为
\[
	\rho^\vee(g)=\rho(g^{-1})^{\mathsf T}=\rho(g)^{-\mathsf T}.
\]
直接转置会使乘法次序反转：\(\rho(gh)^{\mathsf T}=\rho(h)^{\mathsf T}\rho(g)^{\mathsf T}\)，所以当作用矩阵不交换时，它不满足左表示的乘法律。取逆也反转次序，两次反转使 \(\rho^\vee(gh)=\rho^\vee(g)\rho^\vee(h)\)；同时
\[
	\bigl\langle\rho^\vee(g)\varphi,\rho(g)v\bigr\rangle =\varphi\bigl(\rho(g)^{-1}\rho(g)v\bigr) =\langle\varphi,v\rangle.
\]
这里 \(\langle\varphi,v\rangle=\varphi(v)\) 是求值配对。

\ProseParagraphBreak
前面二维余模的加法参数作用，在基 \(v_1,v_2\) 及其对偶基下写成
\[
	\rho(c)=\begin{pmatrix}1&c\\0&1\end{pmatrix},\qquad
	\rho^\vee(c)=\rho(-c)^{\mathsf T}
	=\begin{pmatrix}1&0\\-c&1\end{pmatrix}
	\quad(c\in k).
\]
由于这些矩阵彼此交换，\(c\mapsto\rho(c)^{\mathsf T}\) 在这个例子中仍是加法群的表示。但直接转置未必保持求值配对：取 \(\varphi=v_1^*\)、\(v=v_2\)，有
\[
	\bigl\langle\rho(c)^{\mathsf T}v_1^*,\rho(c)v_2\bigr\rangle=2c, \qquad \langle v_1^*,v_2\rangle=0.
\]
在特征不为二且 \(c\ne0\) 时，两值不同；特征二时本例的转置恰与逆转置相同。坐标函数的对极 \(S(t)=-t\) 把参数换为相反数，从而给出上述保持配对的逆转置。对矩阵的多项式条目求导，又得到
\[
	N=\left.\frac{\mathrm{d}\rho(c)}{\mathrm{d}c}\right|_{c=0}
	=\begin{pmatrix}0&1\\0&0\end{pmatrix},\qquad
	N^\vee=\left.\frac{\mathrm{d}\rho^\vee(c)}{\mathrm{d}c}\right|_{c=0} =-N^{\mathsf T}.
\]
这也恢复了\cref{b5:prop:lie-dual-tensor-hom}的对偶作用：包络代数中 \(S(x)=-x\)，所以 \((x\varphi)(v)=-\varphi(xv)\)。群的逆转置和 Lie 代数的负转置，都是同一对极公式在生成元上的表现。

\ProseParagraphBreak
普通交换 \(v\otimes w\mapsto w\otimes v\) 是否等变，要看余乘法是否余交换；\subsecref{b5:subsec:B01:03}给出一个非余交换的例子。上述基础可对照\cite[Definitions 5.2.2, 5.3.2; Propositions 5.3.5--5.3.6; Corollary 5.3.7]{EtingofEtAl2015TensorCategories}。

% !TeX root = ../../Book5.tex
% 纲要标题：### 21.2 代数群的坐标环
% 纲要位置：工作记录/计划纲要.md，第 4764 行。

\section{代数群的坐标环}
\label{b5:sec:2102}


% 纲要细目（保留原定写作范围）：
%
% * 仿射代数群
% * 群乘法与 Hopf 结构
% * 一般线性群及可对角化群
% * 特殊线性群到射影一般线性群的中心商与表示下降
% * 用微分展示判别中心商的平展性，以双数值点说明基域核平凡不足以保证表示下降
%

代数群的乘法也可由函数上的拉回来描述。采用坐标环而非只看一个基域上的点，有两个好处：公式可以同时在不同交换代数上取值，正特征中的幂零信息也不会因为点集太小而消失。

\subsection{仿射代数群}
\label{b5:subsec:2102:01}

\begin{definition}\label{b5:def:affine-group}
	设 \(k\) 为域，\(A\) 为 \(k\) 上的有限生成交换 Hopf 代数，结构映射记为 \(\Delta,\varepsilon,S\)。对每个交换含幺 \(k\) 代数 \(R\)，记 \(\eta_R:k\to R\) 为单位映射，并令
	\[
		G(R)=\operatorname{Hom}_{k\text{-alg}}(A,R),
	\]
	对 \(a\in A\)，用 Sweedler 记号写 \(\Delta(a)=\sum a_{(1)}\otimes a_{(2)}\)。对 \(g,h\in G(R)\)，用余乘法、余单位及对极分别定义
	\[
		(gh)(a)=\sum g\bigl(a_{(1)}\bigr)h\bigl(a_{(2)}\bigr),\qquad 1_G=\eta_R\varepsilon,\qquad g^{-1}=g\circ S.
	\]
	这个群值函子称为一个有限型仿射群概形，记为 \(G=\operatorname{Spec}A\)，其坐标环记为 \(k[G]=A\)。在代数闭域上，若坐标环约化，也把相应的仿射代数集及其群运算称为\term[fangshe-daishuqun]{仿射代数群}[affine algebraic group]。
\end{definition}
\symidx{\(G(R),k[G]\)：仿射群概形的代数值点群与坐标环}
因为 \(A,R\) 都交换，两个代数同态的卷积仍为代数同态；Hopf 公理正好使以上运算结合、有单位与逆。一个 \(k\) 代数同态 \(R\to R'\) 与这些公式交换，所以这是对 \(R\) 自然的群，而非彼此不相干的一组点集。
\ProseParagraphBreak
两个这类群之间的态射，定义为反向的 Hopf 代数同态。特别地，点上的映射必须在所有 \(R\) 上由同一坐标公式给出；只给 \(G(k)\to H(k)\) 的一个抽象群同态，并不足以定义代数群态射。

\subsection{群乘法与 Hopf 结构}
\label{b5:subsec:2102:02}

坐标环方法与通常的正则映射方法相符。设 \(k\) 代数闭，\(G\) 为仿射代数集，乘法、单位及求逆都是正则映射。则乘法拉回给出
\[
	\Delta:k[G]\to k[G\times G]=k[G]\otimes k[G],
\]
单位与求逆分别给出 \(\varepsilon,S\)。正则函数在三元组 \((g,h,\ell)\) 上比较时，群乘法结合律就是余结合律，其他公理同理。反向从约化的有限型交换 Hopf 代数出发，正则函数关系恢复同样的运算。
\begin{proposition}\label{b5:prop:hopf-ideals-subgroups}
	设 \(k\) 为域，\(A=k[G]\) 为 \(k\) 上的有限生成交换 Hopf 代数，其结构映射为 \(\Delta,\varepsilon,S\)。若理想 \(I\subseteq A\) 满足
	\[
		\Delta(I)\subseteq I\otimes A+A\otimes I,\qquad \varepsilon(I)=0,\qquad S(I)\subseteq I,
	\]
	则 \(A/I\) 自然成为交换 Hopf 代数，且 \(\operatorname{Spec}(A/I)\) 是 \(G\) 的闭子群概形。反过来，使商上的三种映射良定义所需的条件正是以上三式。
\end{proposition}
\begin{proof}
	张量商映射 \(A\otimes A\to(A/I)\otimes(A/I)\) 的核为 \(I\otimes A+A\otimes I\)，故第一条件使余乘法下降；其余两条件分别使余单位与对极下降。所有公理在商中继承。对每个 \(R\)，商到 \(R\) 的代数同态就是那些零化 \(I\) 的 \(G(R)\) 元素，以上三式保证它们对乘法、单位和逆封闭。
\end{proof}
例如 \(k[t]\) 上的 \(\Delta(t)=t\otimes1+1\otimes t\) 描述加法群 \(\mathbb G_a\)。特征 \(p\) 中 \(I=(t^p)\) 满足上述条件，因为二项式中间系数为零；所得闭子群概形有非零幂零坐标，不能仅由其 \(k\) 点辨认。
\symidx{\(\mathbb G_a\)：加法群概形}

\subsection{一般线性群及可对角化群}
\label{b5:subsec:2102:03}

一般线性群的坐标代数为
\[
	k[\operatorname{GL}_n] =k[x_{ij},d]/\bigl(d\det(x_{ij})-1\bigr).
\]
对任意 \(R\)，它的 \(R\) 点就是可逆矩阵。矩阵乘法给出
\[
	\Delta(x_{ij})=\sum_\ell x_{i\ell}\otimes x_{\ell j}, \qquad \varepsilon(x_{ij})=\delta_{ij},
\]
对极则是逆矩阵的条目，由伴随矩阵与行列式逆表示。公式 \(\det(AB)=\det A\det B\) 保证 \(d\) 的余乘法为 \(d\otimes d\)。关系 \(\det X=1\) 又给闭子群 \(\operatorname{SL}_n\)。
\begin{definition}\label{b5:def:diagonalizable-group}
	设 \(k\) 为域，\(M\) 为有限生成交换群。群代数 \(k[M]\) 以符号 \(e^m\) 为基，乘法为 \(e^me^n=e^{m+n}\)，余乘法为 \(\Delta(e^m)=e^m\otimes e^m\)。相应的交换仿射群概形记为 \(D(M)\)，称为分裂\term[keduijiaohua-qun]{可对角化群}[diagonalizable group]。当 \(M=\mathbb Z^r\) 时，它是分裂\term[daishu-huanmian]{代数环面}[algebraic torus] \(\mathbb G_m^r\)。
\end{definition}
\symidx{\(D(M),\mathbb G_m\)：分裂可对角化群与乘法群}
事实上 \(D(M)(R)=\operatorname{Hom}_{\mathrm{grp}}(M,R^\times)\)，因为代数同态必须把 \(e^m\) 送到相容的单位。特别地，
\[
	k[\mathbb G_m]=k[t,t^{-1}],\qquad D(\mathbb Z/n\mathbb Z)=\mu_n,\quad k[\mu_n]=k[t]/(t^n-1).
\]
若 \(k\) 特征为 \(p\)，则 \(t^p-1=(t-1)^p\)，\(\mu_p\) 的坐标环不约化；即使 \(k\) 代数闭，它的 \(k\) 点也只有单位。与此相对，常数群 \(C_p\) 的坐标环为 \(k^{C_p}\)，约化且有 \(p\) 个 \(k\) 点。这是两个不同群概形，不能只因为都带下标 \(p\) 就认同。
\symidx{\(\mu_n\)：\(n\) 次单位根群概形}

\ProseParagraphBreak
可逆矩阵的射影类组成\term[sheying-yiban-xianxingqun]{射影一般线性群}[projective general linear group] \(\operatorname{PGL}_n\)。作为概形，它是射影矩阵空间 \(\mathbb P_k^{n^2-1}\) 中行列式不为零的仿射开子概形 \(D_+(\det)\)，乘法由矩阵乘法给出，逆由伴随矩阵给出。 在任意交换 \(k\) 代数 \(R\) 上，射影类可在开覆盖上用可逆矩阵表示，重叠处的代表相差一个单位标量；当 \(R\) 为域时，每个类都有全局矩阵代表。 下面的局部计算说明，特殊线性群到这个群的商映射也有坐标环层面的含义。
\begin{proposition}\label{b5:prop:special-projective-quotient}
	设 \(k\) 为域，\(n\ge2\)。取矩阵的射影类给出群概形态射
	\[
		q:\operatorname{SL}_n\longrightarrow\operatorname{PGL}_n.
	\]
	该态射有限且忠实平坦，核为中心子群概形 \(\mu_n\)，其元素以标量矩阵嵌入。对任意有限维 \(k\) 空间 \(V\)，一个群概形态射 \(r:\operatorname{SL}_n\to\operatorname{GL}(V)\) 唯一地因子化为某个群概形态射 \(\bar r:\operatorname{PGL}_n\to\operatorname{GL}(V)\) 与 \(q\) 的复合，当且仅当 \(r\) 在 \(\mu_n\) 上平凡。 因此 \(\operatorname{PGL}_n\) 是此中心商，记为 \(\operatorname{SL}_n/\mu_n\)。
\end{proposition}
\begin{proof}
	在 \(\operatorname{PGL}_n\) 的一个标准仿射开集 \(U=\operatorname{Spec}B\) 上，将某个非零射影坐标归一化为一，得到通用矩阵 \(M\)，其行列式 \(D\) 在 \(B\) 中可逆。 \(q^{-1}(U)\) 上的矩阵恰为 \(tM\)，条件是 \(t^nD=1\)，故其坐标环为
	\[
		C=B[t]/(t^n-D^{-1}).
	\]
	它是以 \(1,t,\ldots,t^{n-1}\) 为基的自由 \(B\) 模，而且 \(t\) 可逆。 这些开集覆盖目标，所以 \(q\) 有限且忠实平坦；特别地，它在概形上满射。 单位射影类的原像是 \(tI\)、\(t^n=1\)，即 \(\mu_n\)。 在上述开集上，\(C\otimes_BC\) 中的两个标量参数 \(t,u\) 满足同一个 \(n\) 次幂方程，故 \(z=u/t\) 满足 \(z^n=1\)。这给出
	\[
		C\otimes_BC\cong C[z]/(z^n-1),\qquad u=tz.
	\]
	从右向左的环同态送 \(z\) 到 \(u/t\)；反向保持第一份 \(C\)，并送第二份的 \(u\) 到 \(tz\)。因为 \((tz)^n=t^n\)，第二份的定义关系得到保持；两次复合又分别恢复 \(u\) 与 \(z\)，所以这确为互逆同构。 这些同构在重叠处相容，因而 \(\operatorname{SL}_n\times\mu_n\cong \operatorname{SL}_n\times_{\operatorname{PGL}_n}\operatorname{SL}_n\)， 其中两投影分别为 \((g,z)\mapsto g\) 和 \((g,z)\mapsto gz\)。

	还须验证因子化保持正则性，而不只是在点群上取商。 在上述坐标环中，\(\mu_n\) 的作用对应
	\[
		C\longrightarrow C\otimes_k k[z]/(z^n-1),\qquad t\longmapsto t\otimes z,\quad b\longmapsto b\otimes1.
	\]
	若 \(c=\sum_{i=0}^{n-1}b_it^i\) 在此作用下不变，比较 \(1,z,\ldots,z^{n-1}\) 的系数便知 \(b_i=0\) 对所有 \(i>0\) 成立。 因此不变子环恰为 \(B\)，这一计算不要求 \(n\) 在 \(k\) 中可逆。 当 \(r\) 在 \(\mu_n\) 上平凡时，\(r\) 的矩阵条目及行列式逆在该作用下不变，因而都是 \(B\) 中的函数。 它们在 \(U\) 上定义到 \(\operatorname{GL}(V)\) 的态射。 重叠处的两组函数拉回后相同，而这些自由扩张是忠实平坦的，故原函数也相同；于是局部态射拼成唯一的 \(\bar r\)。 在 \(q\times q\) 下拉回后，\(\bar r\) 的乘法相容性就是 \(r\) 的乘法相容性，忠实平坦性再使该等式在目标上成立。 故 \(\bar r\) 是群概形态射。反向若已有因子化，则 \(\ker q=\mu_n\) 显然作用平凡。
\end{proof}
% 跨卷引用目标：b3:thm:universal-differentials（14.1.3）；b3:def:finite-presentation-infinitesimal-properties（14.5.2）。
回顾第三卷定理14.1.3的微分语言：对交换环同态 \(B\to C\)，\term[kahler-weifen-mo]{Kähler 微分模}[module of Kähler differentials]\footnote{凯勒（Erich Kähler，1906-2000），德国数学家，研究复几何，并把微分方法用于代数几何。} \(\Omega_{C/B}\) 记录所有取值于 \(C\) 模的 \(B\) 导子。按第三卷定义14.5.2，一个有限展示交换 \(B\) 代数 \(C\)，若作为 \(B\) 模平坦且 \(\Omega_{C/B}=0\)，则相应的仿射概形态射称为\term[pingzhan]{平展}[étale]。有限自由性只保证平坦性，还须检查微分模是否消失。
\symidx{\(\Omega_{C/B}\)：交换代数的相对 Kähler 微分模}
\begin{corollary}\label{b5:cor:special-projective-etale}
	设 \(k\) 为域，\(n\ge2\)。中心商态射 \(q:\operatorname{SL}_n\to\operatorname{PGL}_n\) 平展，当且仅当 \(n\) 在 \(k\) 中可逆。若 \(\operatorname{char}k=p>0\) 且 \(p\mid n\)，这个态射有限忠实平坦，但不平展。
\end{corollary}
\begin{proof}
	% 跨卷引用目标：b3:cor:differential-presentation（14.2.2）。
	在\cref{b5:prop:special-projective-quotient}的仿射开集上，仍记 \(C=B[t]/(t^n-D^{-1})\)、\(D\in B^\times\)。这个代数有限展示且作为 \(B\) 模自由。第三卷推论14.2.2的微分展示，把唯一方程的微分作为唯一关系，给出
	\[
		\Omega_{C/B}=C\,\mathrm{d}t/\bigl(nt^{n-1}\cdot\mathrm{d}t\bigr).
	\]
	因为 \(t\) 可逆，若 \(n\in k^\times\)，关系系数 \(nt^{n-1}\) 是单位，故微分模为零，各开集上的态射平展。若 \(\operatorname{char}k=p\mid n\)，则 \(n\cdot1_k=0\)，该微分模为非零的自由模 \(C\,\mathrm{d}t\)，所以态射不平展；原证明的自由基仍然成立，有限忠实平坦性不受影响。
\end{proof}
特别地，当 \(n=p=\operatorname{char}k\) 时，在任意代数闭扩域 \(K/k\) 的几何点处，纤维的坐标环形如 \(K[t]/(t^p-a)\)，其中 \(a\in K^\times\)。选 \(b\in K\) 满足 \(b^p=a\)，便得到 \(K[t]/((t-b)^p)\)。纤维只有一个几何点，却带有非零幂零元；单位纤维正是非约化的 \(\mu_p\)。
\begin{proposition}\label{b5:prop:standard-representation-nondescent}
	设 \(k\) 为特征 \(p>0\) 的域。特殊线性群 \(\operatorname{SL}_p\) 在 \(k^p\) 上的标准有理表示，在 \(\mu_p(k)\) 上平凡，但不能下降为 \(\operatorname{PGL}_p\) 的有理表示。
\end{proposition}
\begin{proof}
	域中 \(z^p=1\) 等价于 \((z-1)^p=0\)，故 \(\mu_p(k)=\{1\}\)。然而在双数环 \(R=k[\epsilon]/(\epsilon^2)\) 中，\((1+\epsilon)^p=1\)，所以
	\[
		z=(1+\epsilon)I_p\in\operatorname{SL}_p(R)
	\]
	是 \(\mu_p(R)\) 的元素，在 \(\operatorname{PGL}_p(R)\) 中成为单位。标准表示仍把它送到 \((1+\epsilon)I_p\ne I_p\)，因而在核群概形 \(\mu_p\) 上不平凡。由\cref{b5:prop:special-projective-quotient}的下降判据，该表示不能下降。
\end{proof}
若再假设 \(k\) 代数闭，则 \(\operatorname{SL}_p(k)\to\operatorname{PGL}_p(k)\) 反而是抽象群同构。它的核只有单位；对任意可逆矩阵 \(A\)，取 \(\lambda\in k^\times\) 满足 \(\lambda^p=(\det A)^{-1}\)，矩阵 \(\lambda A\) 就是同一射影类的行列式一代表。因此可以通过这个点群同构转移标准作用，却不能得到有理群概形表示。表示下降所要求的是核在所有交换 \(k\) 代数上的作用平凡，基域点上的检查不足以保证这一条件。

% !TeX root = ../../Book5.tex
% 纲要标题：### 21.3 有理表示
% 纲要位置：工作记录/计划纲要.md，第 4772 行。

\section{有理表示}
\label{b5:sec:2103}


% 纲要细目（保留原定写作范围）：
%
% * 余模与有理表示
% * 权和代数环面
% * 与 Lie 代数表示的关系及积分问题
% * 用矩阵系数恢复加法群坐标环，比较可对角化群与常数群的模特征表示，并计算矩阵 Lie 括号及 Frobenius 拉回
%

代数群表示不能只要求每个群元素给出一个矩阵，还要让矩阵条目随群元素正则地变化。余模的系数恰好记录这些函数，因此能把“有理表示”的条件写成有限个代数等式。这里的有理表示是正则群态射，不是允许任意带极点的有理映射。

\subsection{余模与有理表示}
\label{b5:subsec:2103:01}

\begin{definition}\label{b5:def:rational-representation}
设 \(G=\operatorname{Spec}A\) 为域 \(k\) 上的有限型仿射群概形，\(V\) 为有限维 \(k\) 空间。群概形态射 \(G\to\operatorname{GL}(V)\) 称为一个有限维\term[youli-biaoshi]{有理表示}[rational representation]。
\end{definition}
\begin{theorem}\label{b5:thm:rational-comodule}
设 \(k\) 为域，\((A,\Delta,\varepsilon,S)\) 为 \(k\) 上的有限生成交换 Hopf 代数，\(G=\operatorname{Spec}A\)。有限维有理 \(G\) 表示与有限维右 \(A\) 余模等价。具体地，给定有限维右余模 \((V,\delta)\) 及交换含幺 \(k\) 代数 \(R\)，若在 \(V\) 的一组 \(k\) 基 \(v_1,\ldots,v_n\) 下有
\[
\delta(v_j)=\sum_i v_i\otimes a_{ij},
\]
则对应的 \(R\) 点 \(g:A\to R\) 在 \(V\otimes_k R\) 上的作用矩阵为 \(\bigl(g(a_{ij})\bigr)\)，且
\[
\Delta(a_{ij})=\sum_\ell a_{i\ell}\otimes a_{\ell j},\qquad
\varepsilon(a_{ij})=\delta_{ij}.
\]
\end{theorem}
\begin{proof}
余模的两条公理分别给出所列矩阵系数等式。对 \(R\) 点求值，前式表示乘法相容，后式表示单位作用为恒等。对极给出逆矩阵 \(\bigl(S(a_{ij})\bigr)\)，因此这些矩阵在每个 \(R\) 上都可逆，定义到 \(\operatorname{GL}_n\) 的态射。反过来，群态射的矩阵系数必满足同样等式，定义余作用。换基只是同时改变这组系数的表达，同态相容条件也相同，所以得到基无关的范畴等价。
\end{proof}
\begin{lemma}\label{b5:lem:comodule-locally-finite}
设 \(k\) 为域，\(C\) 为 \(k\) 余代数，\(V\) 为右 \(C\) 余模，不要求有限维。每个 \(v\in V\) 都属于一个有限维 \(k\) 子余模。
\end{lemma}
\begin{proof}
写 \(\delta(v)=\sum_{i=1}^r v_i\otimes c_i\)，令 \(c_i\) 线性无关。选线性泛函 \(\phi_i\) 满足 \(\phi_i(c_j)=\delta_{ij}\)。令
\(W=\operatorname{span}\{v_i\}\)，亦即所有 \((1\otimes\phi)\delta(v)\) 的线性生成空间。
余结合律在最后一个因子上应用任意 \(\phi\)，给出
\[
\delta\bigl((1\otimes\phi)\delta(v)\bigr)
=\sum_i v_i\otimes(1\otimes\phi)\Delta(c_i)\in W\otimes C.
\]
故 \(W\) 稳定；余单位公式使 \(v=\sum_i\varepsilon(c_i)v_i\in W\)。
\end{proof}
\begin{corollary}\label{b5:thm:affine-group-linear}
域 \(k\) 上每个有限型仿射群概形都可以闭嵌入某个 \(\operatorname{GL}_n\)。
\end{corollary}
\begin{proof}
将 \(A=k[G]\) 本身看作余作用为 \(\Delta\) 的右余模。取 \(A\) 的有限个代数生成元，利用前一引理，选包含这些元及一的有限维子余模 \(W\)。在 \(W\) 的基 \(w_j\) 上写系数为 \(a_{ij}\)。这里使用 \(A\) 作为余代数的另一侧余单位等式 \((\varepsilon\otimes1)\Delta=1_A\)，而非右余模的 \((1\otimes\varepsilon)\Delta=1_A\)。由
\(w_j=(\varepsilon\otimes1)\Delta(w_j)=\sum_i\varepsilon(w_i)a_{ij}\)，
所有 \(w_j\) 都在这些矩阵系数的线性生成空间中。因此表示给出的
\(k\bigl[\operatorname{GL}(W)\bigr]\to A\) 的像包含全部代数生成元，是满射；这就是闭嵌入。
\end{proof}
这项证明同时说明有限生成坐标环的用途。余模局部有限性保证任何一个函数都受有限维矩阵系数控制，而有限生成性保证一次选择便能控制整个坐标环。
\ProseParagraphBreak
在加法群 \(\mathbb G_a\) 的坐标环 \(A=k[t]\) 中，可以直接取 \(W=\operatorname{span}_k\{1,t\}\)。因为
\[
\Delta(1)=1\otimes1,\qquad
\Delta(t)=t\otimes1+1\otimes t,
\]
这个空间是右子余模；在依次排列的基 \(1,t\) 下，矩阵系数为
\[
\begin{pmatrix}1&t\\0&1\end{pmatrix}.
\]
因此所得表示 \(\rho:\mathbb G_a\to\operatorname{GL}_2\) 的坐标同态把 \(x_{12}\) 送到 \(t\)，把 \(x_{11},x_{22}\) 送到 \(1\)，把 \(x_{21}\) 送到 \(0\)，并把行列式逆送到 \(1\)。像已包含 \(k[t]\) 的代数生成元 \(t\)，故这个同态满射，\(\rho\) 为闭嵌入。这是这个二维余模的又一个用途：矩阵条目不仅给点上的作用，还能恢复整个坐标环。

\subsection{权和代数环面}
\label{b5:subsec:2103:02}

\begin{theorem}\label{b5:thm:torus-weight-decomposition}
设 \(k\) 为任意特征的域，\(M\) 为有限生成交换群，\(G=D(M)=\operatorname{Spec}k[M]\)，\(e^m\) 为群代数 \(k[M]\) 中元素 \(m\in M\) 对应的基元。对有限维有理 \(G\) 表示，以及更一般的任意右 \(k[M]\) 余模 \((V,\delta)\)，令 \(V_m=\{v\in V\mid\delta(v)=v\otimes e^m\}\)。则有唯一的分解
\[
V=\bigoplus_{m\in M}V_m,\qquad
\delta(v)=v\otimes e^m\quad(v\in V_m).
\]
有限维时只有有限多个非零项。因此每个有理表示完全可约，简单表示是一维权；结论在任意特征下成立。
\end{theorem}
\begin{proof}
群代数的 \(e^m\) 是一组基，对每个 \(v\) 唯一写
\(\delta(v)=\sum_m v_m\otimes e^m\)，有限多个系数非零。余结合比较第三因子的 \(e^m\) 系数，得到
\(\delta(v_m)=v_m\otimes e^m\)；余单位使 \(v=\sum_m v_m\)。
若一组不同权向量的和为零，对其余作用比较不同 \(e^m\) 系数便知每项为零，所以是直和。表示同态保各权空间，一维子空间在同一权空间内都稳定，故分解为一维表示。唯一性也来自系数唯一性。
\end{proof}
代数群同态 \(\chi:G\to\mathbb G_m\) 称为\term[daishuqun-tezhengbiao]{代数群特征标}[character of an algebraic group]，它给出一维有理表示。这里的特征标是乘法型态射；前面有限群表示的特征标则是表示矩阵的迹。对 \(\mathbb G_m^r\)，权是整数元组 \((m_1,\ldots,m_r)\)，作用特征标为
\(t_1^{m_1}\cdots t_r^{m_r}\)。这不是任意复指数。比如一维交换 Lie 代数的任意复标量作用都给出 Lie 代数表示，却只有整数标量才能从 \(\mathbb G_m\) 的有理表示微分得到。
\ProseParagraphBreak
设 \(\operatorname{char}k=p>0\)。即使 \(G=\mu_p\) 的坐标环 \(k[t]/((t-1)^p)\) 不约化，上述结论仍给出
\[
V=\bigoplus_{j\in\mathbb Z/p\mathbb Z}V_j,\qquad
\delta(v)=v\otimes t^j\quad(v\in V_j).
\]
这里 \(1,t,\ldots,t^{p-1}\) 仍是坐标环的一组 \(k\) 基；权分解的证明用的是这组基及其余乘法，不要求底层代数没有幂零元。因此每个有限维有理 \(\mu_p\) 表示完全可约，尽管该群概形非约化且不光滑。\termidx[guanghua]{光滑}
\ProseParagraphBreak
与此相对，常数群概形 \(C_p\) 的坐标环 \(k^{C_p}\) 是 \(p\) 个域 \(k\) 的直积；每个概形分量都是 \(\operatorname{Spec}k\)，所以它约化且光滑，但有理表示不一定完全可约。例如生成元可在 \(kv_1\oplus kv_2\) 上作用为
\[
J=\begin{pmatrix}1&1\\0&1\end{pmatrix}=I+N.
\]
由于 \(N^2=0\)，在特征 \(p\) 中有 \(J^p=I\)，所以这确实给出 \(C_p\) 的表示。直线 \(kv_1\) 稳定；任何补直线可由 \(v_2+av_1\) 生成，而 \(J(v_2+av_1)=v_2+(a+1)v_1\) 不在这条直线上，故没有稳定补空间。
\ProseParagraphBreak
坐标环约化、群概形光滑和有理表示完全可约，分别是不同条件。后面定义的约化群还要求光滑连通，并排除非平凡的光滑连通正规幺幂闭子群；\(\mu_p\) 不光滑，\(C_p\) 不连通，不能把这两个群直接列作那一定义中的约化群。

\subsection{与 Lie 代数表示的关系及积分问题}
\label{b5:subsec:2103:03}

把一个 \(R\) 点在单位附近代入双数环 \(k[\epsilon]/(\epsilon^2)\)，便得到无穷小作用。这样的点只记录平方为零的一阶变化。
\begin{definition}\label{b5:def:group-scheme-lie-algebra}
设 \(k\) 为域，\(G=\operatorname{Spec}A\) 为 \(k\) 上的有限型仿射群概形，\(\varepsilon:A\to k\) 为坐标 Hopf 代数的余单位。满足
\[
\mathrm{d}(ab)=\varepsilon(a)\cdot\mathrm{d}(b)+\mathrm{d}(a)\cdot\varepsilon(b)
\qquad(a,b\in A)
\]
的 \(k\) 线性映射 \(\mathrm{d}:A\to k\) 称为单位处的导子，它们组成 \(G\) 的 Lie 代数 \(\operatorname{Lie}(G)\) 的底向量空间。映射 \(\varepsilon+\epsilon\cdot\mathrm{d}:A\to k[\epsilon]/(\epsilon^2)\) 是代数同态，其模 \(\epsilon\) 的像为单位点；反过来，每个这样的双数值点唯一给出一个 \(\mathrm{d}\)。
\end{definition}
\symidx{\(\operatorname{Lie}(G)\)：仿射群概形单位处导子组成的 Lie 代数}
对两线性映射 \(\mathrm{d}_1,\mathrm{d}_2:A\to k\)，卷积为
\[
(\mathrm{d}_1*\mathrm{d}_2)(a)=\sum \mathrm{d}_1(a_{(1)})\cdot\mathrm{d}_2(a_{(2)}),
\qquad \Delta(a)=\sum a_{(1)}\otimes a_{(2)}.
\]
若 \(\mathrm{d}_1,\mathrm{d}_2\) 都满足上述导子等式，利用 \(\Delta\) 的乘法相容性与余单位等式，展开得到
\[
\begin{aligned}
(\mathrm{d}_1*\mathrm{d}_2)(ab)
&=\varepsilon(a)\cdot(\mathrm{d}_1*\mathrm{d}_2)(b)+(\mathrm{d}_1*\mathrm{d}_2)(a)\cdot\varepsilon(b)\\
&\quad+\mathrm{d}_1(a)\cdot\mathrm{d}_2(b)+\mathrm{d}_2(a)\cdot\mathrm{d}_1(b).
\end{aligned}
\]
交换 \(\mathrm{d}_1,\mathrm{d}_2\) 后，最后两项不变，因而卷积交换子 \([\mathrm{d}_1,\mathrm{d}_2]=\mathrm{d}_1*\mathrm{d}_2-\mathrm{d}_2*\mathrm{d}_1\) 仍满足导子等式。卷积的结合律使这个括号满足交替性与 Jacobi 恒等式，所以确实定义了 \(\operatorname{Lie}(G)\) 的 Lie 代数结构。

\ProseParagraphBreak
有理表示的余作用诱导 \(\rho_*(\mathrm{d})=(1\otimes\mathrm{d})\delta\)。余结合律给出
\[
\begin{aligned}
\rho_*(\mathrm{d}_1)\rho_*(\mathrm{d}_2)
&=(1\otimes\mathrm{d}_1\otimes\mathrm{d}_2)(\delta\otimes1)\delta\\
&=(1\otimes\mathrm{d}_1\otimes\mathrm{d}_2)(1\otimes\Delta)\delta\\
&=(1\otimes(\mathrm{d}_1*\mathrm{d}_2))\delta.
\end{aligned}
\]
两边交换 \(\mathrm{d}_1,\mathrm{d}_2\) 后相减，便有
\[
[\rho_*(\mathrm{d}_1),\rho_*(\mathrm{d}_2)]=\rho_*([\mathrm{d}_1,\mathrm{d}_2]).
\]
因此这是一个 Lie 代数表示。下面把卷积括号还原为矩阵交换子，并计算行列式方程给切空间施加的条件。
\begin{proposition}\label{b5:prop:matrix-group-lie-bracket}
设 \(k\) 为任意特征的域，\(n\ge1\)。单位处导子与双数值矩阵 \(I+\epsilon X\) 的对应给出 Lie 代数同构
\[
\operatorname{Lie}(\operatorname{GL}_n)\cong\mathfrak{gl}_n(k),\qquad
\operatorname{Lie}(\operatorname{SL}_n)\cong\mathfrak{sl}_n(k)
=\{X\in M_n(k)\mid\operatorname{tr}X=0\},
\]
两者的括号都对应 \([X,Y]=XY-YX\)。
\end{proposition}
\begin{proof}
在双数环中，任何模 \(\epsilon\) 为 \(I\) 的矩阵唯一写成 \(I+\epsilon X\)，其逆为 \(I-\epsilon X\)。于是每个 \(X\in M_n(k)\) 都给出 \(\operatorname{GL}_n\) 的单位处导子 \(\mathrm{d}_X\)，满足 \(\mathrm{d}_X(x_{ij})=X_{ij}\)，并且每个单位处导子都如此得到。由矩阵坐标的余乘法，
\[
\begin{aligned}
(\mathrm{d}_X*\mathrm{d}_Y)(x_{ij})
&=\sum_{\ell=1}^n\mathrm{d}_X(x_{i\ell})\cdot\mathrm{d}_Y(x_{\ell j})\\
&=\sum_{\ell=1}^n X_{i\ell}Y_{\ell j}=(XY)_{ij}.
\end{aligned}
\]
所以 \([\mathrm{d}_X,\mathrm{d}_Y](x_{ij})=(XY-YX)_{ij}\)。这决定了整个坐标环上的导子：矩阵条目生成多项式部分，行列式逆的值由逆元求导法则确定。因此卷积交换子正对应矩阵交换子。

行列式的置换展开中，模 \(\epsilon^2\) 后，只有恒等置换的各个一次对角项留下，得到
\[
\det(I+\epsilon X)=1+\epsilon\operatorname{tr}X.
\]
故这个点属于 \(\operatorname{SL}_n\)，当且仅当 \(\operatorname{tr}X=0\)。对闭子群限制坐标函数后，导子的卷积仍使用同一余乘法，所以括号也是矩阵交换子；\(\operatorname{tr}(XY-YX)=0\) 又表明这个子空间对括号封闭。
\end{proof}
特别地，在特征 \(p\) 中，标量矩阵 \(I_p\) 的迹为零。\(\mu_p\) 的双数值点 \((1+a\epsilon)I_p\) 给出切方向 \(aI_p\)，所以 \(\operatorname{Lie}(\mu_p)=kI_p\) 非零。中心商 \(q:\operatorname{SL}_p\to\operatorname{PGL}_p\) 将所有这些点送到单位，其微分便零化这个标量切方向。前面标准表示不能下降的例子，正是保留了点集上看不到的这一方向。
\begin{proposition}\label{b5:prop:additive-rational-representations}
设 \(k\) 为特征零域，\(V\) 为 \(n\) 维 \(k\) 向量空间。加法群 \(\mathbb G_a\) 在 \(V\) 上的每个有理表示，在选定 \(k\) 基后唯一形如
\[
\rho(t)=\exp(tN)=\sum_{j\ge0}\frac{t^jN^j}{j!},
\]
其中 \(N\) 为幂零矩阵，故该和有限。每个有限维有理 \(\mathbb G_m\) 表示的微分则可对角化，特征值为整数。
\end{proposition}
\begin{proof}
加法群的矩阵条目为 \(t\) 的多项式，满足
\(\rho(t+u)=\rho(t)\rho(u)\)、\(\rho(0)=I\)。
对 \(u\) 求导并令 \(u=0\)，得到 \(\rho'(t)=\rho(t)N\)，其中 \(N=\rho'(0)\)。逐次比较多项式系数得到 \(A_j=N^j/j!\)。多项式次数有限，所以充分大的 \(N^j=0\)。反过来，幂零 \(N\) 的有限指数满足加法公式，给出有理表示。乘法群的结论由\cref{b5:thm:torus-weight-decomposition}中 \(t^m\) 的导数得到。
\end{proof}
因此从 Lie 代数表示到代数群有理表示会多出条件；在特征零时，即使两群的 Lie 代数都一维交换，\(\mathbb G_a\) 要求幂零作用，\(\mathbb G_m\) 要求整数半单作用。复 Lie 群的积分问题还涉及拓扑覆盖，留在附录C说明。
\begin{example}\label{b5:exm:frobenius-rational-representation}
设 \(k\) 为特征 \(p>0\) 的域。在 \(k^2\) 上考虑加法群 \(\mathbb G_a\) 的两个有理表示
\[
\rho_0(t)=I_2,\qquad
\rho_1(t)=\begin{pmatrix}1&t^p\\0&1\end{pmatrix}.
\]
恒等式 \((s+t)^p=s^p+t^p\) 对任意交换 \(k\) 代数成立，故 \(\rho_1(s+t)=\rho_1(s)\rho_1(t)\)；矩阵条目为多项式且行列式为一，因而确实定义有理表示。它是前面二维加法群表示与群态射 \(F:t\mapsto t^p\) 的复合。

对双数值参数 \(t=a\epsilon\)，有 \((a\epsilon)^p=0\)，所以 \(\rho_1(a\epsilon)=I_2=\rho_0(a\epsilon)\)。两表示的 Lie 代数作用都为零，却不是同一个有理表示：它们的右上角矩阵系数分别为坐标多项式 \(t^p\) 与 \(0\)。平凡表示的矩阵在任意换基后仍为 \(I_2\)，而 \(\rho_1\) 不是恒等作用，所以这两个有理表示也不同构。一阶微分因而不能恢复这个 Frobenius 拉回所保留的高次项。这里的相同只指单位处的一阶作用；在 \(R=k[u]/(u^{p+1})\) 中取参数 \(t=u\)，便有 \(u^p\ne0\)，两表示在这个含幂零元的试验代数上已经不同。
\end{example}

% !TeX root = ../../Book5.tex
% 纲要标题：### 21.4 代数群结构初步
% 纲要位置：工作记录/计划纲要.md，第 4779 行。

\section{代数群结构初步}
\label{b5:sec:2104}


% 纲要细目（保留原定写作范围）：
%
% * 幺幂群和约化群
% * 根数据的意义
% * 正特征行为和群概形方向
% * 计算三阶幺幂群的中心与 Lie 括号，区别一般线性群的根方向、中心方向及其整数格
%

从矩阵例子进入一般代数群时，需要区分几个不同的结构条件。幺幂性控制元素的 Jordan 形，约化性排除连通正规幺幂部分；根数据又在根系以外保留了整数格。正特征时，坐标环的幂零信息和无穷小子群还会使点集直觉失效。

\subsection{幺幂群和约化群}
\label{b5:subsec:2104:01}

以下先在代数闭域 \(k\) 上讨论光滑仿射代数群，闭子群也取光滑闭子群。矩阵 \(u\) 若满足 \(u-I\) 幂零，就称为幺幂矩阵。一个光滑线性代数群在一个忠实表示中全部 \(k\) 点都幺幂时，称为幺幂群；这个概念的表示无关性是一般线性代数群理论的一项结构结论。本节的实例均在明确的矩阵模型中使用它。
\ProseParagraphBreak
上三角对角线全为一的群 \(U_n\) 是基本例子。令 \(J\) 为上三角矩阵代数中的严格上三角矩阵理想，\(U_n^{(r)}=1+J^r\)，即仅在距对角线至少 \(r\) 条的上对角线上允许非零条目的子群。因为 \(J\) 幂零，\(1+J^{r+s}\) 是 \(1+J\) 的正规子群。在商群中，若 \(X\in J^r\)、\(Y\in J^s\)，则 \(XY,YX\in J^{r+s}\)，所以
\[
	(1+X)(1+Y)\equiv1+X+Y\pmod{J^{r+s}} \quad(X\in J^r,\ Y\in J^s),
\]
与逆序相乘得到的 \((1+Y)(1+X)\) 相同。因此两个元素在商群中交换，其交换子属于 \(1+J^{r+s}\)，即 \(\bigl[U_n^{(r)},U_n^{(s)}\bigr]\subseteq U_n^{(r+s)}\)。因为 \(J^n=0\)，这个群幂零，特别可解。
\begin{proposition}\label{b5:prop:unipotent-three-heisenberg}
	设 \(k\) 为域，\(U_3\) 为 \(k\) 上对角线全为一的三阶上三角矩阵群，记
	\[
		u(a,b,c)=\begin{pmatrix}1&a&c\\0&1&b\\0&0&1\end{pmatrix},
		\qquad a,b,c\in k.
	\]
	则点群 \(U_3(k)\) 的中心与交换子群相同，均为 \(\{u(0,0,c):c\in k\}\)，且
	\[
		[u(a,b,c),u(a',b',c')]=u(0,0,ab'-a'b).
	\]
	其 Lie 代数为 \(\mathfrak u_3=kE_{12}\oplus kE_{23}\oplus kE_{13}\)，括号由
	\[
		[E_{12},E_{23}]=E_{13},\qquad [E_{12},E_{13}]=[E_{23},E_{13}]=0
	\]
	按交替双线性延拓，中心与导出代数都为 \(kE_{13}\)。因此它是三维 Heisenberg Lie 代数，且该点群与 Lie 代数的幂零类都为二。
\end{proposition}
\begin{proof}
	相乘及求逆得到
	\[
		\begin{aligned}
			u(a,b,c)u(a',b',c')&=u(a+a',b+b',c+c'+ab'),\\
			u(a,b,c)^{-1}&=u(-a,-b,ab-c).
		\end{aligned}
	\]
	依次代入这两个式子，在约定 \([u,v]=uvu^{-1}v^{-1}\) 下得到
	\[
		[u(a,b,c),u(a',b',c')]=u(0,0,ab'-a'b).
	\]
	若此式对全部 \(a',b'\) 为单位，分别取 \((a',b')=(0,1)\)、\((1,0)\)，便知 \(a=b=0\)。反过来，\(u(0,0,c)\) 与全部元素交换。因此群中心为 \(\{u(0,0,c):c\in k\}\)。每个中心元素也都是交换子，因为 \([u(1,0,0),u(0,c,0)]=u(0,0,c)\)；故交换子群等于中心，\(U_3(k)\) 的幂零类为二。

	单位处的双数值点为 \(I+\epsilon X\)，其中 \(X=aE_{12}+bE_{23}+cE_{13}\)。由\cref{b5:prop:matrix-group-lie-bracket}中闭矩阵子群的括号计算，\(\mathfrak u_3=\operatorname{Lie}(U_3)\) 的括号为矩阵交换子。于是
	\[
		[aE_{12}+bE_{23}+cE_{13},\, a'E_{12}+b'E_{23}+c'E_{13}] =(ab'-a'b)E_{13}.
	\]
	特别地，\([E_{12},E_{23}]=E_{13}\)，而 \(E_{13}\) 与全部 \(\mathfrak u_3\) 交换；同样取两个基方向检验，Lie 中心恰为 \(kE_{13}\)。把\cref{b5:exm:abelian-heisenberg-lie}中的 \(x,y,z\) 分别送到 \(E_{12},E_{23},E_{13}\)，就得到三维 Heisenberg Lie 代数的这个矩阵模型。群交换子与 Lie 括号在中心方向上具有相同的双线性表达式；在当前二步幂零模型中，三次矩阵乘积均为零。一般幺幂矩阵群的交换子还可能含有高次项，不能据此把所有群乘法都当作加法。
\end{proof}
\begin{definition}\label{b5:def:reductive-group}
	设 \(G\) 为代数闭域上的光滑连通仿射代数群。若 \(G\) 没有非平凡的光滑连通正规幺幂闭子群，就称为\term[yuehua-qun]{约化群}[reductive group]。若此外没有正维数的光滑连通正规可解闭子群，就称为半单代数群。这里的“约化群”不同于“坐标环约化”：后者只要求没有非零幂零元。
\end{definition}
例如环面没有非平凡幺幂元素，所以约化；它自身交换且有正维数时不半单。\(U_n\) 在 \(n\ge2\) 时光滑连通且非平凡，却是幺幂的，因而不约化。一般约化群的结构还涉及环面、根子群和正规子群；下一小节从整数格资料观察其中一部分。

\subsection{根数据的意义}
\label{b5:subsec:2104:02}

\begin{definition}\label{b5:def:root-datum}
	设 \(X,Y\) 为两个有限秩自由交换群，给定完美整数配对 \(\langle\ ,\ \rangle:X\times Y\to\mathbb Z\)，以及有限子集 \(\Phi\subseteq X,\Phi^\vee\subseteq Y\) 及双射 \(\alpha\mapsto\alpha^\vee\)。要求对所有 \(\alpha\in\Phi\)，都有 \(\langle\alpha,\alpha^\vee\rangle=2\)，且反射
	\[
		s_\alpha(x)=x-\langle x,\alpha^\vee\rangle\alpha,\qquad s_\alpha^\vee(y)=y-\langle\alpha,y\rangle\alpha^\vee
	\]
	分别保持 \(\Phi\) 和 \(\Phi^\vee\)，并对所有 \(\alpha,\beta\in\Phi\) 满足 \((s_\alpha\beta)^\vee=s_\alpha^\vee(\beta^\vee)\)。还要求约化条件 \(\Phi\cap\mathbb Q\alpha=\{\alpha,-\alpha\}\) 对每个 \(\alpha\in\Phi\) 成立，其中把 \(X\) 嵌入 \(X\otimes_{\mathbb Z}\mathbb Q\)。满足这些条件的资料称为\term[gen-shuju]{根数据}[root datum]；允许 \(\Phi\) 为空，也不要求它生成整个 \(X\otimes\mathbb Q\)，以便包含环面的中心方向。
\end{definition}
\symidx{\(X,Y,\Phi,\Phi^\vee\)：根数据中的特征标格、余特征标格、根与余根集合}
在连通约化群的结构理论中，若 \(T\) 是极大环面，这两个格具体为
\[
	X=\operatorname{Hom}_{\mathrm{alg.grp}}(T,\mathbb G_m),\qquad Y=\operatorname{Hom}_{\mathrm{alg.grp}}(\mathbb G_m,T).
\]
这里的同态是代数群态射，\(\chi\in X\) 称为特征标，\(\nu\in Y\) 称为一参数子群。它们的复合是 \(\mathbb G_m\) 的群态射，故形如 \(t\mapsto t^n\)，其唯一整数指数就是配对：
\[
	\chi(\nu(t))=t^{\langle\chi,\nu\rangle}.
\]
这一整数幂结论也可直接由 Laurent 多项式的 Hopf 等式看出：类群元只能是单项式 \(t^n\)。根系给出实向量空间里的角度与相对长度，根数据还记录哪些整数权给出极大环面上的特征标。这些特征标不必延拓为整个群的一维表示；例如 \(\operatorname{SL}_2\) 的环面有非平凡特征标，而 \(\operatorname{SL}_2\) 自身没有非平凡的一维有理表示。因此同一个复半单 Lie 代数可以对应不同的代数群。
\ProseParagraphBreak
例如在代数闭域 \(k\) 上，对 \(G=\operatorname{GL}_2\) 取标准对角环面 \(T\)。记特征标 \(\varepsilon_i\) 为第 \(i\) 个对角条目，即
\[
	\varepsilon_i\bigl(\operatorname{diag}(t_1,t_2)\bigr)=t_i, \qquad X(T)=\mathbb Z\varepsilon_1\oplus\mathbb Z\varepsilon_2.
\]
对角矩阵共轭 \(E_{12}\)、\(E_{21}\) 的系数分别为 \(t_1/t_2\)、\(t_2/t_1\)，对角矩阵本身则被固定。因此非零根只有 \(\pm\alpha\)、\(\alpha=\varepsilon_1-\varepsilon_2\)，只张成一维有理子空间。行列式特征标 \(\delta=\varepsilon_1+\varepsilon_2\) 给出另一个方向：在中心环面 \(Z(G)=\{tI_2\}\) 上，\(\alpha(tI_2)=1\)，而 \(\delta(tI_2)=t^2\)。所以
\[
	X(T)\otimes_{\mathbb Z}\mathbb Q =\mathbb Q\alpha\oplus\mathbb Q\delta,
\]
其中行列式方向记录根所看不到的中心环面。
\ProseParagraphBreak
这里的直和是在有理数上成立的，\(\alpha,\delta\) 并非特征标格的整数基。因为 \(\varepsilon_1=\dfrac{\delta+\alpha}{2}\)、\(\varepsilon_2=\dfrac{\delta-\alpha}{2}\)， \(\mathbb Z\alpha+\mathbb Z\delta\) 只包含两个坐标同奇偶的特征标，在 \(X(T)\) 中指数为二。完整特征标格仍须保留 \(\varepsilon_1,\varepsilon_2\) 所给的整数资料。若 \(\operatorname{char}k\ne2\)，在 Lie 代数层面也可写 \(\mathfrak{gl}_2(k)=\mathfrak{sl}_2(k)\oplus kI_2\)；这个分解来自把矩阵减去其迹的一半乘单位矩阵，中心方向同样不参与非零根空间。
\ProseParagraphBreak
对 \(\operatorname{SL}_2\)，取 \(T=\bigl\{\operatorname{diag}(t,t^{-1})\bigr\}\)。若 \(\lambda(t)=t\)，则 \(X=\mathbb Z\lambda\)，根为 \(\pm2\lambda\)；一参数子群 \(\nu(t)=\operatorname{diag}(t,t^{-1})\) 满足 \(\langle\lambda,\nu\rangle=1\)，正余根就是 \(\nu\)。 对 \(\operatorname{PGL}_2\) 的对角环面，用对角条目比值 \(r\) 作坐标，特征标格由 \(\alpha(r)=r\) 生成，根为 \(\pm\alpha\)。若 \(\eta(t)\) 是 \(\operatorname{diag}(t,1)\) 的射影类，则 \(\langle\alpha,\eta\rangle=1\)。商映射把 \(\operatorname{SL}_2\) 环面的参数送为 \(r=t^2\)，所以特征标拉回为 \(\alpha\mapsto2\lambda\)，正余根 \(\nu\) 的像则为 \(2\eta\)。于是 \(\operatorname{PGL}_2\) 的余根是格 \(Y\) 的生成元的二倍，仍满足根余根配对为二。两者的抽象根系都为 \(A_1\)，但整数格中的嵌入不同。
\ProseParagraphBreak
这也解释了表示下降的一个必要条件。中心元 \(-I\in\operatorname{SL}_2(\mathbb C)\) 在 \(L_m=\operatorname{Sym}^m(\mathbb C^2)\) 上作用为 \((-1)^m\)。由\cref{b5:prop:special-projective-quotient}，代数群商 \(\operatorname{SL}_{2,\mathbb C}/\mu_2\) 同构于 \(\operatorname{PGL}_{2,\mathbb C}\)，它在复点上的商群也记为 \(\operatorname{PSL}_2(\mathbb C)\)。所以这个有理表示能下降到该商代数群当且仅当 \(m\) 为偶数。Lie 代数的最高权分类列出所有 \(m\ge0\)，却不会自动替每个中心商筛去不能下降的表示。

\subsection{正特征下的代数群与群概形}
\label{b5:subsec:2104:03}

在特征 \(p\) 中，对加法群的 Frobenius 态射
\[
	F:\mathbb G_a\longrightarrow\mathbb G_a,\qquad t\longmapsto t^p
\]
对应坐标同态 \(k[t]\to k[t]\)、\(t\mapsto t^p\)。它的核为 \(\alpha_p=\operatorname{Spec}k[t]/(t^p)\)。
\symidx{\(\alpha_p\)：特征 \(p\) 中加法群 Frobenius 态射的核}
若 \(k\) 代数闭，点集上 \(F\) 为双射，但坐标同态不满，所以它不是代数群同构。其微分为零；核概形的切空间却一维，尽管它的 \(k\) 点只有一个。
\ProseParagraphBreak
同样，\(\mu_p=\operatorname{Spec}k[t,t^{-1}]/(t^p-1)\) 的切空间一维。对任意域扩张 \(K/k\)，方程 \(t^p=1\) 等价于 \((t-1)^p=0\)；域中没有非零幂零元，所以 \(\mu_p(K)=\{1\}\)。但在双数环 \(R=k[\epsilon]/(\epsilon^2)\) 中，对每个 \(a\in k\) 都有
\[
	(1+a\epsilon)^p=1+a^p\epsilon^p=1.
\]
其中 \(1+a\epsilon\) 可逆，逆为 \(1-a\epsilon\)，因此确定一个 \(\mu_p(R)\) 点；若 \(a\ne0\)，这个点就不是单位点。只看域值点会遗漏这些无穷小变化；允许含幂零元的代数作为试验环，才能在点函子中保留它们。
\ProseParagraphBreak
Frobenius 态射与非约化核说明，光滑性、约化性、可分性以及表示是否能由微分恢复，分别控制不同信息。后续的不变量理论将区分特征零的线性约化性与正特征下更弱的几何约化性；有限群的平均公式只在群阶可逆时适用，一般群概形需要相应的表示论方法。

% !TeX root = ../../Book5.tex
% 纲要标题：### 21.5 型、对合与经典群
% 纲要位置：工作记录/计划纲要.md，第 4786 行。
% 纲要细目（保留原定写作范围）：
%
% * 从非退化型的分裂矩阵模型解释正交、辛、酉群及相似变换群，写清底域、特征和型的约定
% * 区别等距群、相似变换群、代数自同构群与中心商；代数群满态射不保证有理点映射满射
% * 选一个四元数的非分裂模型，并重述所需代数事实；不以一般扭曲形式分类或未写的非交换上同调作为前提
% * 以行列式限制乘子，计算二维辛型的三个群，比较四元数的实 Lie 型、复化及三维旋转商
%

\section{型、对合与经典群}
\label{b5:sec:2105}

% 纲要小节：
% #### 21.5.1 等距群、相似变换群与乘子
% #### 21.5.2 带对合代数的自同构群
% #### 21.5.3 分裂模型与四元数例子
% ---

正交群、辛群和酉群最初由保持型的矩阵定义。把型换成自同态代数上的伴随对合以后，同一批方程也能在非分裂代数中使用。本节先固定具体模型，再比较等距、相似变换和代数自同构；它们之间有联系，却不是同一个群。

\subsection{等距群、相似变换群与乘子}
\label{b5:subsec:2105:01}

设 \(F\) 为特征不为二的域，取 \(K=F\) 或一个可分二次扩张 \(K/F\)，并以 \(a\mapsto\bar a\) 分别表示恒等或非平凡共轭。给定 \(H\in\operatorname{GL}_n(K)\)、\(n\ge1\)，假设 \(H^*=\epsilon H\)、\(\epsilon\in\{1,-1\}\)，其中 \(X^*=\bar X^{\mathsf T}\)。型为 \(b(u,v)=\bar u^{\mathsf T}Hv\)，伴随对合为 \(\sigma_H(X)=H^{-1}X^*H\)。
\begin{definition}\label{b5:def:classical-similitude}
	设 \(F\) 为特征不为二的域，\(K=F\) 或 \(K/F\) 为可分二次扩张，共轭分别为恒等或非平凡域自同构。取 \(H\in\operatorname{GL}_n(K)\)、\(n\ge1\)，满足 \(H^*=\epsilon H\)、\(\epsilon\in\{1,-1\}\)，其中 \(X^*=\bar X^{\mathsf T}\)。对非退化型 \(b(u,v)=\bar u^{\mathsf T}Hv\)，定义等距群与相似变换群
	\[
		\begin{aligned}
			\operatorname{Isom}(b)&=\bigl\{g\in\operatorname{GL}_n(K)\mid g^*Hg=H\bigr\},\\
			\operatorname{Sim}(b)&=\bigl\{g\in\operatorname{GL}_n(K)\mid g^*Hg=cH\;\text{对某个}\;c\in F^\times\bigr\}.
		\end{aligned}
	\]
	相似变换的这个唯一标量 \(c\) 称为乘子，记为 \(\mu(g)\)。当 \(K=F\)、\(H\) 对称时，两群分别记为 \(\operatorname O(b)\)、\(\operatorname{GO}(b)\)；当 \(H\) 交替时，记为 \(\operatorname{Sp}(b)\)、\(\operatorname{GSp}(b)\)；二次扩张上的 Hermitian 情形记为 \(\operatorname U(b)\)、\(\operatorname{GU}(b)\)。
\end{definition}
\symidx{\(\operatorname{Isom}(b),\operatorname{Sim}(b),\mu(g)\)：型的等距群、相似变换群与乘子}
\symidx{\(\operatorname O(b),\operatorname{Sp}(b),\operatorname U(b)\)：正交、辛与酉等距群}
\symidx{\(\operatorname{GO}(b),\operatorname{GSp}(b),\operatorname{GU}(b)\)：正交、辛与酉相似变换群}
这些式子不仅定义 \(F\) 上的抽象群。对任意交换 \(F\) 代数 \(R\)，记 \(K_R=K\otimes_FR\)，并把共轭延拓为固定 \(R\) 的自同构。把 \(H\) 的系数映入 \(K_R\) 后，两群在 \(R\) 上的点由
\[
	\begin{aligned}
		\operatorname{Isom}(b)(R) &=\{g\in\operatorname{GL}_n(K_R):g^*Hg=H\},\\
		\operatorname{Sim}(b)(R) &=\{(g,c)\in\operatorname{GL}_n(K_R)\times R^\times:g^*Hg=cH\}
	\end{aligned}
\]
给出。乘子唯一，所以第二式与前面的矩阵定义相容。选取 \(K\) 的一个 \(F\) 基，可把这些式子写成有限个 \(F\) 坐标方程，得到仿射 \(F\) 群概形。它们的环境群函子是 \(R\mapsto\operatorname{GL}_n(K_R)\) 及其与 \(\mathbb G_m\) 的积；当 \(K=F\) 时，前者就是 \(\operatorname{GL}_{n,F}\)。

\ProseParagraphBreak

二次扩张上若 \(H^*=-H\)，可取 \(u\in K\) 满足 \(u\ne\bar u\)，令 \(c=u-\bar u\)。则 \(c\ne0\)、\(\bar c=-c\)，从而 \((cH)^*=cH\)。把型乘以 \(c\) 不改变等距群、相似变换群及乘子，也不改变伴随对合 \(\sigma_H\)。因此这种型的两群也可用相应 Hermitian 型的酉群记号 \(\operatorname U\)、\(\operatorname{GU}\) 表示。
\begin{proposition}\label{b5:prop:classical-multiplier}
	设 \(F\) 为特征不为二的域，\(K=F\) 或 \(K/F\) 为可分二次扩张，共轭分别为恒等或非平凡域自同构，\(X^*=\bar X^{\mathsf T}\)。取 \(n\ge1\)、\(H\in\operatorname{GL}_n(K)\)，满足 \(H^*=\epsilon H\)、\(\epsilon=\pm1\)。令 \(b\) 为 \(K^n\) 上的非退化型 \(b(u,v)=\bar u^{\mathsf T}Hv\)，\(\sigma_H(X)=H^{-1}X^*H\)，记 \(\operatorname{Sim}(b)\)、\(\operatorname{Isom}(b)\) 分别为相似变换群和等距群。则乘子给出群同态 \(\mu:\operatorname{Sim}(b)\to F^\times\)，其核为 \(\operatorname{Isom}(b)\)。并且
	\[
		g\in\operatorname{Sim}(b),\ \mu(g)=c \quad\Longleftrightarrow\quad \sigma_H(g)g=cI.
	\]
	每个相似变换还满足 \(\overline{\det g}\det g=\mu(g)^n\)。
\end{proposition}
\begin{proof}
	非退化型至少有一个非零配对值，故乘子唯一。两个相似变换连续作用使配对值乘两个标量的积，逆变换乘以逆标量，所以是同态，乘子一恰是等距条件。最后左乘 \(H^{-1}\) 即将矩阵方程改写为所示对合方程。

	对 \(g^*Hg=cH\) 取行列式，得到
	\[
		\overline{\det g}\det(H)\det g=c^n\det(H).
	\]
	因为 \(\det H\ne0\)，消去它就得到 \(\overline{\det g}\det g=c^n\)。
\end{proof}
这个等式在二次扩张情形为 \(N_{K/F}(\det g)=c^n\)，而在 \(K=F\) 时为 \((\det g)^2=c^n\)，给出乘子的一个必要条件。例如 \(F=K=\mathbb R\)、\(n\) 为奇数时，右端必须为正，所以每个非退化实对称型都只能有正乘子，不需要假定型正定。但行列式条件未必充分：正定实型即使在偶数维也不允许负乘子，因为对非零向量 \(v\)，等式 \(b(gv,gv)=c\,b(v,v)\) 中的两个型值 \(b(gv,gv)\)、\(b(v,v)\) 都为正。
\ProseParagraphBreak
乘子的实际像还须结合型判断。在正定实型上，\(\sqrt c I\) 实现每个正数；在分裂辛型
\(J=\begin{psmallmatrix}0&I_m\\-I_m&0\end{psmallmatrix}\)
上，\(\operatorname{diag}(cI_m,I_m)\) 实现每个 \(c\in F^\times\)。所以写核与像时不能一概补上“到 \(F^\times\) 满射”。

\subsection{带对合代数的自同构群}
\label{b5:subsec:2105:02}

\begin{proposition}\label{b5:prop:involution-automorphism-similitude}
	设 \(F\) 为特征不为二的域，\(K=F\) 或 \(K/F\) 为可分二次扩张，共轭分别为恒等或非平凡域自同构，\(X^*=\bar X^{\mathsf T}\)。取 \(H\in\operatorname{GL}_n(K)\)、\(n\ge1\)，满足 \(H^*=\epsilon H\)、\(\epsilon=\pm1\)。令 \(A=M_n(K)\)、\(b(u,v)=\bar u^{\mathsf T}Hv\) 为 \(K^n\) 上的非退化型、\(\sigma(X)=H^{-1}X^*H\)。以下比较所给域上的点群：\(\operatorname{Sim}(b)\) 指由 \(g\in\operatorname{GL}_n(K)\) 组成的 \(F\) 有理点相似变换群，保持 \(\sigma\) 的 \(K\) 代数自同构组成的群记为 \(\operatorname{Aut}_K(A,\sigma)\)。内自同构 \(\operatorname{Int}(g)(X)=gXg^{-1}\) 给出群同构
	\[
		\operatorname{Sim}(b)/K^\times \ \cong\ \operatorname{Aut}_K(A,\sigma),
	\]
	其中 \(K^\times\) 作为非零标量矩阵嵌入相似群，其乘子为 \(\bar a a\)。
\end{proposition}
\begin{proof}
	矩阵代数的每个 \(K\) 代数自同构都是内自同构，这是第二卷定理17.3.1的 Skolem--Noether 结论\footnote{斯科伦（Thoralf Albert Skolem，1887-1963），挪威数学家，研究数理逻辑与集合论；他关于一阶逻辑和可数模型的工作是其代表成果。}在中心单代数 \(M_n(K)\) 上的应用。直接计算
	\[
		\sigma(gXg^{-1})=\sigma(g)^{-1}\sigma(X)\sigma(g).
	\]
	它等于 \(g\sigma(X)g^{-1}\) 对所有 \(X\) 成立，当且仅当 \(\sigma(g)g\) 与所有矩阵交换，即为某个 \(cI\)、\(c\in K^\times\)。对该等式再用 \(\sigma\)，得到 \(\bar c=c\)，故 \(c\in F^\times\)。由\cref{b5:prop:classical-multiplier}，这正是相似条件。内自同构的核是矩阵代数的中心单位 \(K^\times\)，而这些标量全属于相似群，遂得所示同构。
\end{proof}
自同构要求逐点固定 \(K\)。二次扩张上若允许自同构在中心上也作共轭，会得到另一个群。对一般带对合中心单代数，相同的计算仍把内自同构与条件 \(\sigma(g)g\in F^\times\) 联系起来；非分裂情形中的型及对合细节见第二卷附录B。
\ProseParagraphBreak
在特征不为二的域 \(F\) 上，二维交替型把三个点群同时化为熟悉的矩阵群。取 \(K=F\)、\(J=\begin{psmallmatrix}0&1\\-1&0\end{psmallmatrix}\)，对任意 \(g=\begin{psmallmatrix}a&b\\c&d\end{psmallmatrix}\in\operatorname{GL}_2(F)\)，直接相乘得到
\[
	g^{\mathsf T}Jg
	=\begin{pmatrix}0&ad-bc\\bc-ad&0\end{pmatrix}
	=(\det g)J.
\]
因此 \(\operatorname{Sim}(J)=\operatorname{GL}_2(F)\)，其乘子就是行列式；乘子核给出 \(\operatorname{Isom}(J)=\operatorname{SL}_2(F)\)。再由\cref{b5:prop:involution-automorphism-similitude}，对 \(\sigma_J(X)=J^{-1}X^{\mathsf T}J\) 有
\[
	\operatorname{Aut}_F(M_2(F),\sigma_J) \cong\operatorname{GL}_2(F)/F^\times I_2 =\operatorname{PGL}_2(F).
\]
这里从相似群到对合自同构群的映射是内共轭，所去掉的正是非零标量矩阵。
\ProseParagraphBreak
还须区分代数群满态射与域上点的满射。由\cref{b5:prop:special-projective-quotient}， \(\operatorname{SL}_{2,\mathbb R}\to\operatorname{PGL}_{2,\mathbb R}\) 是核为 \(\mu_2\) 的有限忠实平坦商态射，因而在群概形层面满射。复数上，标量平方根能把任意可逆二阶矩阵缩放为行列式一矩阵，所以复点映射也满射。实点的像却只含行列式为正的射影矩阵类：缩放把行列式乘一个正平方，不可能把负行列式改成一。对二维辛型，\(g=\operatorname{diag}(-1,1)\) 的乘子为 \(-1\)，故它是相似变换而非等距变换；其内共轭仍保持 \(\sigma_J\)，但其射影类没有实的 \(\operatorname{SL}_2\) 原像。保持对合的自同构在这个模型中都来自相似变换，却未必来自等距变换。

\subsection{分裂模型与四元数例子}
\label{b5:subsec:2105:03}

在 \(F=\mathbb R\) 上，Hamilton 四元数代数
\[
	\mathbb H=\mathbb R\oplus\mathbb Ri\oplus\mathbb Rj\oplus\mathbb Rk, \quad i^2=j^2=k^2=-1,\quad ij=k=-ji
\]
带有共轭 \(\bar q\) 和范数 \(N(q)=q\bar q=q_0^2+q_1^2+q_2^2+q_3^2\)。 共轭反转乘法，故 \(N(uv)=N(u)N(v)\)。于是
\[
	\operatorname{SL}_1(\mathbb H)(\mathbb R)=\bigl\{q\in\mathbb H\mid N(q)=1\bigr\}
\]
由一个四元二次方程定义为实代数群；逆由共轭给出，不需要除以变量。
\ProseParagraphBreak
复化后它变成分裂模型。把 \(i,j\) 分别送到
\[
	\begin{pmatrix}\mathrm i&0\\0&-\mathrm i\end{pmatrix},
	\qquad
	\begin{pmatrix}0&1\\-1&0\end{pmatrix}.
\]
两矩阵平方为 \(-I\)、反交换，并且它们连同单位及乘积线性无关，所以给出 \(\mathbb H\otimes_{\mathbb R}\mathbb C\cong M_2(\mathbb C)\)。 计算一般元素的行列式恰为 \(N(q)\)，因此复化后的群为 \(\operatorname{SL}_2\)。在实点上，四元数单位球面是有界闭集，而 \(\operatorname{SL}_2(\mathbb R)\) 含任意大的 \(\operatorname{diag}(t,t^{-1})\)；这两个实模型的差异不由其共同的复化消除。
\ProseParagraphBreak
单位处的切空间也可以直接从范数方程求出。记 \(G=\operatorname{SL}_1(\mathbb H)\)，在双数环 \(\mathbb R[\epsilon]/(\epsilon^2)\) 中有
\[
	N(1+\epsilon u) =(1+\epsilon u)(1+\epsilon\bar u) =1+\epsilon(u+\bar u).
\]
所以 \(1+\epsilon u\) 是 \(G\) 在单位处的双数值点，当且仅当 \(u+\bar u=0\)，即 \(\operatorname{Lie}(G)=\operatorname{Im}\mathbb H=\mathbb Ri\oplus\mathbb Rj\oplus\mathbb Rk\)。四元数左乘给出 \(G\) 在 \(\mathbb H\) 上的忠实有理表示，其微分把 \(u\) 送到左乘算子 \(L_u\)。因为 \(L_u(1)=u\)，这个微分是单射。由\cref{b5:prop:matrix-group-lie-bracket}及微分的括号相容性，\(L_uL_v=L_{uv}\) 使括号成为 \([u,v]=uv-vu\)，故
\[
	[i,j]=2k,\qquad [j,k]=2i,\qquad [k,i]=2j.
\]
上述复矩阵模型把 \(i,j,k\) 送到三个线性无关的迹零矩阵，因而也给出 \(\operatorname{Lie}(G)\otimes_{\mathbb R}\mathbb C\cong\mathfrak{sl}_2(\mathbb C)\)。
\ProseParagraphBreak
这个实 Lie 代数与 \(\mathfrak{sl}_2(\mathbb R)\) 仍有区别。在纯虚空间的基 \(i,j,k\) 下取欧氏内积 \(\langle\ ,\ \rangle\) 和叉积 \(\times\)，乘法表给出 \(uv=-\langle u,v\rangle+u\times v\)，因此 \(\operatorname{ad}_u(v)=2u\times v\)。于是
\[
	(\operatorname{ad}_u)^2(v) =4\bigl(\langle u,v\rangle u-N(u)v\bigr),\qquad \operatorname{tr}\bigl((\operatorname{ad}_u)^2\bigr)=-8N(u).
\]
这对每个非零 \(u\) 都为负。相反，在 \(\mathfrak{sl}_2(\mathbb R)\) 中，\(h=\operatorname{diag}(1,-1)\) 的伴随作用在 \(h,E_{12},E_{21}\) 上的特征值为 \(0,2,-2\)，所以 \(\operatorname{tr}((\operatorname{ad}_h)^2)=8>0\)。Lie 同构把对应的伴随算子共轭，必须保持这一平方迹，故这两个实 Lie 代数不同构。
\begin{example}\label{b5:exm:quaternion-conjugation-group}
	单位四元数通过 \(C_q(x)=qxq^{-1}\) 作用于纯虚三维实空间 \(\operatorname{Im}\mathbb H\)，且 \(C_qC_r=C_{qr}\)。因为 \(q^{-1}=\bar q\)，对 \(\bar x=-x\) 有
	\[
		\overline{qx\bar q}=q\bar x\bar q=-qx\bar q, \qquad N(qx\bar q)=N(x),
	\]
	所以 \(C_q\) 保持纯虚空间及其正定范数。作用的核是 \(\{\pm1\}\)：若与 \(i,j,k\) 都交换，便在 \(\mathbb H\) 的中心 \(\mathbb R\) 中，再加范数一就只剩两个元。具体地，与 \(i\) 交换先使 \(j,k\) 的系数为零，再与 \(j\) 交换使 \(i\) 的系数为零。
	\ProseParagraphBreak
	取单位纯虚四元数 \(u\) 及实数 \(\theta\)，令
	\[
		q=\cos\frac\theta2+u\sin\frac\theta2.
	\]
	它的范数为一。对 \(x=au+y\)、\(\langle u,y\rangle=0\)，有 \(uy=-yu=u\times y\)、\(uyu=y\)。展开 \(qxq^{-1}\)，得到
	\[
		C_q(au+y)=au+\cos\theta\,y+\sin\theta\,(u\times y).
	\]
	因此它固定 \(u\)，并在 \(u^\perp\) 上作角度为 \(\theta\) 的旋转。每个非实的单位四元数都可如此表示：把纯虚部分归一化为 \(u\)，其实部与纯虚部分的长度满足平方和为一；实的单位四元数 \(\pm1\) 则作用为恒等。故全部 \(C_q\) 都属于 \(\operatorname{SO}_3(\mathbb R)\)。
	\ProseParagraphBreak
	反过来，对 \(R\in\operatorname{SO}_3(\mathbb R)\)，非实特征值成共轭对且模为一，实特征值为 \(\pm1\)。由于维数为三、行列式为一，它至少有一个特征值一，因而固定某个单位向量 \(u\)。在 \(u^\perp\) 上的限制是行列式一的二维正交变换，故在基 \(v,u\times v\) 下形如一个平面旋转矩阵，其中 \(v\) 为任意单位向量且 \(v\perp u\)。取其旋转角 \(\theta\)，上述四元数 \(q\) 就实现 \(R\)。于是共轭作用满射到 \(\operatorname{SO}_3(\mathbb R)\)，结合已经算出的核，得到点群同构
	\[
		\operatorname{SL}_1(\mathbb H)(\mathbb R)/\{\pm1\} \cong\operatorname{SO}_3(\mathbb R).
	\]
	例如 \(q=i\) 固定 \(i\) 并把 \(j,k\) 变号；而 \(q=(1+i)/\sqrt2\) 固定 \(i\)，将 \(j\) 送到 \(k\)、\(k\) 送到 \(-j\)。这是由非交换代数乘法产生的具体正交作用。
\end{example}
分裂矩阵型、非分裂四元数、等距群与代数自同构因而可以互相解释。它们也提醒我们：复数域上的根系或 Lie 代数类型只记录部分结构；改换基域、对合或中心商以后，点群与表示下降问题还需要重新核对。


\begin{chaptersummary}
	余乘法规定张量积的作用，余单位规定平凡对象，对极使对偶和求逆成为可能。群代数中的元素按对角方式作用，包络代数中的 Lie 代数生成元按 Leibniz 方式作用；坐标代数则以反向的函数拉回描述群乘法。

	\ProseParagraphBreak
	有理表示等同于坐标余模，有限生成坐标环可由有限维矩阵系数恢复。环面表示按整数权分解，而微分只记录这些特征标在单位处的一阶变化；正特征的 Frobenius 和非约化核说明一阶变化可能遗漏更多信息。

	\ProseParagraphBreak
	型的等距群是相似群的乘子核，保持伴随对合的内自同构则还要模去中心标量。四元数群在复化后成为熟悉的矩阵群，但其实际底域上的点仍可大不相同。这些区别是下一章讨论不变量与商时需要保留的对象条件。
\end{chaptersummary}

\BookExercises
\begin{optionalexercise}\label{b5:ex:21:01}
	设 \((H,\Delta,\varepsilon,S)\) 为域 \(k\) 上的 Hopf 代数，证明 \(\varepsilon\bigl(S(h)\bigr)=\varepsilon(h)\)。提示：在 \(\operatorname{Hom}_k(H,k)\) 的卷积中比较余单位及其逆。
\end{optionalexercise}
\begin{solution}
	对极等式应用 \(\varepsilon\)，得到 \((\varepsilon S)*\varepsilon=\varepsilon=\varepsilon*(\varepsilon S)\)，而余单位公理说明 \(\varepsilon\) 本来就是这个卷积代数的单位。单位的逆仍为单位，故 \(\varepsilon S=\varepsilon\)。
\end{solution}

\begin{optionalexercise}\label{b5:ex:21:02}
	在 \(k[t]\) 上取 \(\Delta(t)=t\otimes1+1\otimes t\)。计算 \(\Delta(t^n)\)，并说明特征 \(p\) 时 \(t^p\) 仍为本原元。特征零时全部本原多项式有哪些？
\end{optionalexercise}
\begin{solution}
	余乘法是代数同态，故 \(\Delta(t^n)=\sum_{j=0}^n\binom njt^j\otimes t^{n-j}\)。特征 \(p\) 下中间二项式系数为零，得到本原性。特征零时若非零本原多项式最高次为 \(d>1\)，其余乘法中 \(t^{d-1}\otimes t\) 系数为首项系数乘 \(d\)，非零，矛盾；常数本原要求 \(c=2c\)，所以为零。全部本原元恰为 \(kt\)。
\end{solution}

\begin{optionalexercise}\label{b5:ex:21:03}
	设 \(G\) 为群。在 \(kG\) 中求满足 \(\Delta(u)=u\otimes u\)、\(\varepsilon(u)=1\) 的全部元素，称这类元素为类群元。
\end{optionalexercise}
\begin{solution}
	写 \(u=\sum a_gg\)。左侧仅有 \(g\otimes g\)，而右侧不同 \(g,h\) 的系数为 \(a_ga_h\)。域中不能有两个非零系数，因此只有一个 \(a_g\ne0\)。比较对角系数得 \(a_g=a_g^2\)，余单位又要求它为一。故类群元恰好是群基 \(G\)。
\end{solution}

\begin{optionalexercise}\label{b5:ex:21:04}
	设 \(G=C_2\)，在特征二中比较 \(kG\) 与 \(k^G\) 的底层代数、类群元及幂零性。
\end{optionalexercise}
\begin{solution}
	前者同构于 \(k[u]/(u^2)\)，其中 \(u=s-1\)，有非零幂零元；后者为 \(k\times k\)，约化。前者类群元为 \(1,s\)。后者的类群元是函数 \(f(gh)=f(g)f(h)\)、\(f(1)=1\)，即特征标 \(G\to k^\times\)；其 \(f(s)^2=1\) 在特征二只有 \(f(s)=1\)，所以只有常数一。两者不能因向量空间维数相同而认同。
\end{solution}

\begin{optionalexercise}\label{b5:ex:21:05}
	在复 \(\mathfrak{sl}_2\) 模 \(L_2\otimes L_2\) 上，令
	\[
		Q=2h\otimes h+4f\otimes e+4e\otimes f.
	\]
	利用 \(\Omega=h^2+2h+4fe\) 的余乘法，求 \(Q\) 的特征值、重数和最小多项式，并写出投到平凡直和项的多项式投影。
\end{optionalexercise}
\begin{solution}
	由\cref{b5:thm:sl2-clebsch-gordan}，有 \(L_2\otimes L_2\cong L_4\oplus L_2\oplus L_0\)。\(\Omega\) 在 \(L_m\) 上的标量为 \(m(m+2)\)，而两个因子上的标量都是八，所以张量积上 \(\Omega=16I+Q\)。因此 \(Q\) 在这三个直和项上依次为标量 \(8,-8,-16\)，重数依次为 \(5,3,1\)。三项均非零且特征值不同，故
	\[
		\mu_Q(t)=(t-8)(t+8)(t+16),\qquad P_0=\dfrac{(Q-8I)(Q+8I)}{192}.
	\]
	\(P_0\) 在特征值八与负八处为零，在负十六处为一，所以其像恰为一维平凡直和项。交叉项把两个因子的作用联系起来，不能由因子上的 Casimir 标量相加代替。
\end{solution}

\begin{optionalexercise}\label{b5:ex:21:06}
	设 \(G\) 为有限群，\(V,W\) 分别是集中在次数 \(a,b\in G\) 的非零 \(G\)-分次 \(k\) 空间，并把它们看作左 \(k^G\) 模。证明普通翻转 \(\tau:V\otimes W\to W\otimes V\) 是 \(k^G\) 模同态当且仅当 \(ab=ba\)。取 \(G=S_3\)、\(a=(12),b=(23)\) 给出不等变的具体例子，并说明为何群代数 \(kG\) 的张量表示没有这一障碍。
\end{optionalexercise}
\begin{solution}
	张量分次公式使 \(V\otimes W\) 集中在 \(ab\) 次，而 \(W\otimes V\) 集中在 \(ba\) 次。若两次相同，翻转保持分次，因而是模同态；若不同，取非零纯张量 \(v\otimes w\)，则
	\[
		\tau\bigl(\delta_{ab}(v\otimes w)\bigr)=w\otimes v, \qquad \delta_{ab}\tau(v\otimes w)=0,
	\]
	故不等变。在 \(S_3\) 的所选两个换位上，\(ab=(123)\)、\(ba=(132)\) 不同；取两空间都一维就得到反例。相反，\(kG\) 的余乘法为 \(\Delta(g)=g\otimes g\)，所以 \(\tau(gv\otimes gw)=gw\otimes gv=g\tau(v\otimes w)\)，不要求 \(G\) 交换。
\end{solution}

\begin{optionalexercise}\label{b5:ex:21:07}
	验证矩阵余代数的公式 \(\Delta(c_{ij})=\sum_\ell c_{i\ell}\otimes c_{\ell j}\)、\(\varepsilon(c_{ij})=\delta_{ij}\)。求其线性对偶代数。
\end{optionalexercise}
\begin{solution}
	两次余乘法均为 \(\sum_{\ell,r}c_{ir}\otimes c_{r\ell}\otimes c_{\ell j}\)，重命名指标即一致；余单位直接只留下一个求和项。若对偶基为 \(e^{ij}\)，卷积满足 \(e^{ij}e^{rs}=\delta_{jr}e^{is}\)，所以对偶代数为矩阵代数。
\end{solution}

\begin{optionalexercise}\label{b5:ex:21:08}
	在 \(\mathbb G_a\) 的正则余模 \(k[t]\) 中，求包含 \(t^d\) 的一个有限维子余模。特征零时证明最小这样的子余模为次数不超过 \(d\) 的多项式空间；特征 \(p\) 时求包含 \(t^p\) 的最小子余模，并与次数不超过 \(p\) 的多项式空间比较。
\end{optionalexercise}
\begin{solution}
	二项式公式说明 \(W_d=\operatorname{span}\{1,t,\ldots,t^d\}\) 稳定。特征零时，令 \(\phi_r\) 取出多项式的 \(t^r\) 系数，则
	\[
		(1\otimes\phi_{d-j})\Delta(t^d)=\binom dj t^j\qquad(0\le j\le d).
	\]
	每个二项式系数非零，故任一含 \(t^d\) 的子余模都含每个 \(t^j\)，最小者就是 \(W_d\)。在特征 \(p\) 中，\(\Delta(t^p)=t^p\otimes1+1\otimes t^p\)，所以 \(W=\operatorname{span}\{1,t^p\}\) 已是子余模。任一含 \(t^p\) 的子余模还须包含 \((1\otimes\phi_p)\Delta(t^p)=1\)，故 \(W\) 最小。它的维数为二，而 \(W_p\) 的维数为 \(p+1\)；中间二项式系数消失，正是最小子余模缩小的原因。
\end{solution}

\begin{optionalexercise}\label{b5:ex:21:09}
	给出 \(\mathbb G_m^2\) 在二维空间上的权 \((1,0),(0,1)\) 作用，计算对偶、二次对称幂和外平方的权。
\end{optionalexercise}
\begin{solution}
	对偶的两个权为 \((-1,0)\)、\((0,-1)\)，因为特征标取逆。对称平方中，基与权的对应为
	\[
		v_1^2:\ (2,0),\qquad v_1v_2:\ (1,1),\qquad v_2^2:\ (0,2).
	\]
	外平方 \(v_1\wedge v_2\) 的权为 \((1,1)\)。这些结论直接由特征标乘法得到，在任意特征均成立。
\end{solution}

\begin{optionalexercise}\label{b5:ex:21:10}
	证明 \(\mathbb G_m\to\mathbb G_m\) 的群概形态射恰为 \(t\mapsto t^n\)、\(n\in\mathbb Z\)。
\end{optionalexercise}
\begin{solution}
	坐标同态由可逆 Laurent 多项式的像 \(u\) 决定，Hopf 条件要求其为类群元。写有限 Laurent 和，比较 \(\Delta(u)\) 与 \(u\otimes u\)，如群代数的系数论证可知 \(u=t^n\)。反向每个这样的公式显然满足全部结构等式。
\end{solution}

\begin{optionalexercise}\label{b5:ex:21:11}
	在特征零域上，证明代数群同态 \(\mathbb G_a\to\mathbb G_m\) 必为平凡，而 \(\mathbb G_a\to\mathbb G_a\) 恰为 \(t\mapsto ct\)。
\end{optionalexercise}
\begin{solution}
	前者对应把 \(k[u,u^{-1}]\) 映到 \(k[t]\)，但 \(k[t]\) 的单位只有非零常数，单位点条件使该常数为一。后者的坐标像必须本原，前题已证全部本原多项式为 \(ct\)，反向加法性显然。
\end{solution}

\begin{optionalexercise}\label{b5:ex:21:12}
	设 \(k\) 特征为零，给 \(\mathbb G_a\) 一个三维有理表示，微分为单个三阶幂零 Jordan 块 \(N\)。写出表示矩阵并求不动空间。
\end{optionalexercise}
\begin{solution}
	以 \(N=\begin{psmallmatrix}0&1&0\\0&0&1\\0&0&0\end{psmallmatrix}\)，有
	\[
		\rho(t)=\begin{pmatrix}1&t&t^2/2\\0&1&t\\0&0&1\end{pmatrix}.
	\]
	一个向量被所有参数固定，当且仅当各个 \(t\) 次系数作用为零，等价于 \(Nv=0\)。因此不动空间是第一坐标轴，维数一。
\end{solution}

\begin{optionalexercise}\label{b5:ex:21:13}
	在特征 \(p\) 中，比较平凡的一维 \(\mathbb G_m\) 表示与特征标 \(t^p\) 的表示：其微分与有理群表示是否相同？
\end{optionalexercise}
\begin{solution}
	平凡特征标的微分为零，\(t^p\) 的导数也为零，所以两个 Lie 代数表示相同。但坐标 Laurent 多项式 \(1\) 与 \(t^p\) 不同，故群概形表示不同；在代数闭域的点上也可选 \(t^p\ne1\)。微分在这里不忠实。
\end{solution}

\begin{optionalexercise}\label{b5:ex:21:14}
	设 \(k\) 代数闭且特征为 \(p\)。计算 \(\alpha_p(k)\)、\(\alpha_p\bigl(k[\epsilon]/(\epsilon^2)\bigr)\)，并与平凡群概形比较。
\end{optionalexercise}
\begin{solution}
	域中 \(a^p=0\) 只可能 \(a=0\)，所以第一点群平凡。双数环中 \(a+b\epsilon\) 的 \(p\) 次幂为 \(a^p\)，故全部 \(b\epsilon\) 都给点，组成与 \((k,+)\) 同构的群。平凡群概形在任何代数上只有一个点，所以两者不同。
\end{solution}

\begin{optionalexercise}\label{b5:ex:21:15}
	设 \(k\) 为域，\(n\ge2\)。证明 \((t^n-1)\) 是 \(k[t,t^{-1}]\) 的 Hopf 理想，求 \(\operatorname{Lie}(\mu_n)\) 及标量嵌入 \(\mu_n\to\operatorname{SL}_n\) 的微分。再用双数值射影矩阵证明
	\[
		\operatorname{Lie}(\operatorname{PGL}_n)\cong\mathfrak{gl}_n/kI,
	\]
	并比较中心商 \(q:\operatorname{SL}_n\to\operatorname{PGL}_n\) 的微分在 \(n\) 可逆及 \(\operatorname{char}k\mid n\) 两种情形下的核与像。
\end{optionalexercise}
\begin{solution}
	有 \(\Delta(t^n-1)=(t^n-1)\otimes t^n+1\otimes(t^n-1)\)，余单位为零，对极将它送为 \(-t^{-n}(t^n-1)\)，故该理想是 Hopf 理想。单位附近的元素为 \(1+\epsilon c\)，方程 \((1+\epsilon c)^n=1+\epsilon nc=1\) 给出 \(\operatorname{Lie}(\mu_n)=\{c\in k:nc=0\}\)。标量嵌入的微分为 \(c\mapsto cI\)；迹为 \(nc\)，所以确实落在 \(\mathfrak{sl}_n\) 中。 双数环 \(R=k[\epsilon]/(\epsilon^2)\) 的单位射影类附近可选矩阵代表 \(I+\epsilon Y\)：它落在任一单位矩阵的非零对角坐标所给的仿射图中，先将该坐标归一化，再用常数标量使模 \(\epsilon\) 的代表为 \(I\)。两个这样的代表相差标量 \(1+\epsilon c\)，即 \(Y\) 相差 \(cI\)。因此切空间为 \(\mathfrak{gl}_n/kI\)。投影 \(\operatorname{GL}_n\to\operatorname{PGL}_n\) 的坐标拉回保持余乘法，故其微分保持导子的卷积及交换子括号；这个微分正是 \(Y\mapsto Y+kI\)，所以所得识别保持 Lie 括号。\(q\) 的微分就是
	\[
		\mathfrak{sl}_n\longrightarrow\mathfrak{gl}_n/kI,\qquad Y\longmapsto Y+kI.
	\]
	若 \(n\) 可逆，每个类可减去 \(\dfrac{\operatorname{tr}Y}{n}I\) 得到唯一迹零代表，故此微分同构。若 \(\operatorname{char}k\mid n\)，则 \(kI\subseteq\mathfrak{sl}_n\)，核为 \(kI\)，像为 \(\mathfrak{sl}_n/kI\)，在 \(\mathfrak{gl}_n/kI\) 中余维一。因此群概形商态射虽满射，其微分在后一情形不满射。
\end{solution}

\begin{optionalexercise}\label{b5:ex:21:16}
	设 \(k\) 为域，\(J\subset M_4(k)\) 为严格上三角矩阵代数，\(U_4(k)=1+J\)。求该点群的下中心列 \(\gamma_1=U_4(k)\)、\(\gamma_{r+1}=[U_4(k),\gamma_r]\)，并与 Lie 代数 \(J\) 的下中心列比较。说明两者的幂零类均为三。
\end{optionalexercise}
\begin{solution}
	正文的交换子估计给出 \(\gamma_2\subseteq1+J^2\)、\(\gamma_3\subseteq1+J^3\)、\(\gamma_4=1\)。记 \(x_{ij}(a)=I+aE_{ij}\)。对 \(i<j<\ell\)，直接相乘得到
	\[
		[x_{ij}(a),x_{j\ell}(b)]=x_{i\ell}(ab).
	\]
	于是全部 \(x_{13}(c),x_{24}(c),x_{14}(c)\) 都属于 \(\gamma_2\)，分别由指标链 \(1<2<3\)、\(2<3<4\)、\(1<2<4\) 实现。它们生成 \(1+J^2\)，故 \(\gamma_2=1+J^2\)。又因 \(x_{24}(c)\in\gamma_2\)，交换子 \([x_{12}(1),x_{24}(c)]=x_{14}(c)\) 给出 \(\gamma_3=1+J^3\ne1\)。因此下中心列为 \(1+J,1+J^2,1+J^3,1\)，幂零类恰为三。 Lie 括号满足 \([J^r,J^s]\subseteq J^{r+s}\)，而
	\[
		[E_{12},E_{23}]=E_{13},\qquad [E_{23},E_{34}]=E_{24},\qquad [E_{12},E_{24}]=E_{14}.
	\]
	这分别生成 \(J^2\) 及非零的 \(J^3\)，故 Lie 代数的下中心列为 \(J,J^2,J^3,0\)，幂零类也为三。这里用具体矩阵计算比较两条列，没有假定一般正特征群与 Lie 代数的结构都能相互恢复。
\end{solution}

\begin{optionalexercise}\label{b5:ex:21:17}
	对实型 \(b\) 的矩阵 \(H=\operatorname{diag}(1,-1)\)，求相似乘子的像，并给出乘子为 \(-1\) 的矩阵。
\end{optionalexercise}
\begin{solution}
	交换坐标的 \(P=\begin{psmallmatrix}0&1\\1&0\end{psmallmatrix}\) 满足 \(P^{\mathsf T}HP=-H\)。正乘子由实标量 \(\sqrt c I\) 实现，负乘子由 \(\sqrt{-c}P\) 实现，故像为全部 \(\mathbb R^\times\)。正定型只能有正乘子，所以型的选择影响实际像。
\end{solution}

\begin{optionalexercise}\label{b5:ex:21:18}
	设 \(b\) 为非退化型，\(a\in K^\times\)。证明标量矩阵 \(aI\) 总给相似变换，但其在自同态代数上的内自同构恒等。由此区分三个群。
\end{optionalexercise}
\begin{solution}
	有 \((aI)^*H(aI)=\bar aaH\)，所以乘子为 \(\bar aa\)。标量与每个矩阵交换，故内自同构恒等。标量矩阵 \(aI\) 属于 \(\operatorname{Isom}(b)\) 当且仅当 \(\bar aa=1\)；相似群允许任意非零标量，而对合代数自同构群已经把这些标量全都商掉。
\end{solution}

\begin{optionalexercise}\label{b5:ex:21:19}
	设 \(F\) 为域，\(n\ge2\)。证明行列式给出满射
	\[
		\operatorname{PGL}_n(F)\longrightarrow F^\times/(F^\times)^n, \qquad [A]\longmapsto\det A\pmod{(F^\times)^n},
	\]
	其核恰为 \(\operatorname{SL}_n(F)\to\operatorname{PGL}_n(F)\) 的像，并求后一个映射的核。用此判定复数及实数域上点映射何时满射。
\end{optionalexercise}
\begin{solution}
	把代表 \(A\) 换成 \(aA\) 时，行列式乘以 \(a^n\)，故映射良定义，且保持乘法。矩阵 \(\operatorname{diag}(c,1,\ldots,1)\) 实现每个 \(c\in F^\times\)，故满射。射影类在核中当且仅当 \(\det A=a^n\) 对某个 \(a\in F^\times\) 成立，此时 \(a^{-1}A\) 行列式一；反向一个行列式一的代表显然使该类在核中。因此
	\[
		\operatorname{PGL}_n(F)/\operatorname{im}\bigl(\operatorname{SL}_n(F)\bigr) \cong F^\times/(F^\times)^n.
	\]
	特殊线性群点映射的核为 \(\{aI:a\in F^\times,\ a^n=1\}\)。复数中每个非零数有 \(n\) 次根，所以点映射总满射。实数中，\(n\) 为奇数时 \(n\) 次幂遍历 \(\mathbb R^\times\)，点映射满射；\(n\) 为偶数时其像恰为正行列式射影类，商为两个元素的符号群。后一情形的核为 \(\{\pm I\}\)，奇数情形的核则只有 \(I\)。
\end{solution}

\begin{optionalexercise}\label{b5:ex:21:20}
	在 \(\mathbb H\) 中令 \(q=(1+i+j+k)/2\)。计算共轭在基 \(i,j,k\) 上的矩阵，以及该矩阵的行列式。
\end{optionalexercise}
\begin{solution}
	四元数范数为一，所以 \(q^{-1}=(1-i-j-k)/2\)。使用 \(ij=k\)、\(jk=i\)、\(ki=j\) 及反向乘积变号，逐项相乘得到 \(qiq^{-1}=j,qjq^{-1}=k,qkq^{-1}=i\)。
	矩阵为 \(\begin{psmallmatrix}0&0&1\\1&0&0\\0&1&0\end{psmallmatrix}\)，是三轮换矩阵，行列式一。它固定 \(i+j+k\)，在其正交平面上是三分之一整周的旋转；\(q\) 与 \(-q\) 仍给相同作用。
\end{solution}

\begin{optionalexercise}\label{b5:ex:21:21}
	对 \(n\ge2,m\ge0\)，证明 \(\operatorname{Sym}^m(\mathbb C^n)\) 的自然有理 \(\operatorname{SL}_{n,\mathbb C}\) 表示能下降为中心商 \(\operatorname{PGL}_{n,\mathbb C}\) 的有理表示，当且仅当 \(n\mid m\)；并在 \(n=2\) 时把条件写成特征标格条件。
\end{optionalexercise}
\begin{solution}
	由\cref{b5:prop:special-projective-quotient}，\(\operatorname{SL}_{n,\mathbb C}\to\operatorname{PGL}_{n,\mathbb C}\) 是核为 \(\mu_n\) 的中心商，有理表示下降当且仅当在该核上平凡。 标量 \(\zeta I\)、\(\zeta^n=1\)，在每个纯 \(m\) 次张量及其对称商上按 \(\zeta^m\) 作用。 复数域上的 \(\mu_n\) 约化，且有本原 \(n\) 次单位根，所以这个子群概形上的作用平凡当且仅当 \(n\mid m\)。 这同时保证了下降表示的正则性，而不只是复点上的群作用。 对 \(n=2\)，环面最高权为 \(m\lambda\)，下降要求它属于根格 \(2\mathbb Z\lambda\)，即 \(m\) 偶数。
\end{solution}

\begin{optionalexercise}\label{b5:ex:21:22}
	取一个交换Hopf代数 \(A\)。证明每个 \(R\) 点的逆由 \(gS\) 给出时，还满足 \((gh)^{-1}=h^{-1}g^{-1}\)，并解释为何不能据此认为 \(G(R)\) 交换。
\end{optionalexercise}
\begin{solution}
	Hopf 对极给卷积逆，任意群中双边逆唯一，故 \((gh)^{-1}=h^{-1}g^{-1}\)。目标代数 \(R\) 的交换性与 \(\Delta\) 保持乘法一起保证卷积结果仍是代数同态；\(A\) 的交换性却不保证 \(\Delta\) 的两个张量因子可以翻转。以 \(A=k[\operatorname{GL}_2]\) 为例，点群仍为通常的非交换矩阵群。
\end{solution}
